% NAL 1644, batch 4 — Gallica views 61–80.
% Pass: Opus 5.5 (claude-opus-5-5), 2026-09-23 — first pass, unchecked against
% the views by a human.
% Read from the local mirror only (archives/nal-1644/f0061.jpg … f0080.jpg):
% a 1000-px thumbnail of each view, then four full-width bands at 2400 px
% (normalised), and tight crops at 1–3× for struck words, marginal notes and
% signs. Each view was drafted as soon as it was read.
%
% Dating. The catalogue gives no dating for this volume: \dating{} is left
% empty, and this pass infers no date.
%
% What the views are. Greyscale scans, portrait, about 3750 × 5620. One page
% per view, rectos at the odd views 61–79, versos at the even views 62–80.
% No view is blank; none is skipped. As in batch 3 the leaves are of very
% unequal size and are bound so that a neighbouring leaf projects above or
% below the one photographed, with its writing: small leaves « 28 », « 31 »,
% « 32 », « 34 », « 35 », « 36 » lie over the larger leaves « 29 », « 33 »
% and « 37 », whose first lines show above them (views 61, 67, 69, 73, 75,
% 77), and the long leaf « 35 » shows its last line below leaves « 33 » and
% « 34 » (views 71–74). Such lines are transcribed only on the view where
% their own page is photographed, and each view's first \note{} says what
% projects and where it belongs. The lines « nempe sunt numeri triangulares
% … quadrans » at the head of leaf « 37 » are thus given at view 79, and
% « metathesi A cubus æquatur B plano in A + Z solido », on a longer leaf
% showing below leaf « 37 » at view 79, belongs to a page after this batch.
%
% What the twenty views hold. One hand throughout, the regular seventeenth-
% century cursive of batch 3, Latin, with its own corrections in the left
% margin (an underlined word or number in the line, the right one in the
% margin between two short strokes). On leaves « 28 »–« 30 » (views 61–66)
% the margin also carries longer notes and calculations in a smaller, tighter
% script — demonstrations of the rules, and the corrected numbers; this pass
% could not tell whether it is a second hand or the same writer writing
% small, and says « la main plus petite » in the notes. Two pieces:
%
%   v. 61–62   Continues Prop. X of « De affectis sub gradibus intermediis »
%              (batch 3, view 60): its exposition and example; Props XI–XIII;
%              how to find three proportional planes whose sum with the mean
%              is a square; « De reliquis », Prop. XIIII.
%   v. 63      Prop. XV (four continued proportionals). Then « Æqualitatum
%              contradicentium constitutiva Cap. XVIII XIX » (the number
%              rewritten), Props I–II.
%   v. 64–65   Props II–V of that chapter (four and six continued
%              proportionals, differences of alternate terms).
%   v. 65–67   « Æqualitatum Inversarum Constitutiva Cap. XIX XX » (XIX
%              struck), Props I–VI.
%   v. 68      « Alia rursus æqualitatum Inversarum Constitutiva Cap. XXI »,
%              Props I–II; the page ends on Prop. II's example, the lower
%              half blank. This is where the chapters of constitutiva stop in
%              this batch.
%   v. 69      A new piece opens at the head of leaf « 32 », in the same hand:
%              « Francisci Vietæ. De æquationum Emendatione tractatus
%              secundus. De solennibus quinque modis præparandarum
%              æquationum adversus earum in numeris δυσμηχανίαν. Ac primum.
%              De expurgatione per uncias quæ remedium est adversus
%              πολυπλοκείαν Cap. I. » — the five modes listed and numbered.
%              A later, small, regular hand has written in the margin « p. 127.
%              editionis 1646. », a reference to a printed edition. The name
%              is the leaf's own; nothing else in the batch is attributed.
%   v. 70–72   Cap. I continued: the unciæ (semis, triens, quadrans…), the
%              three cases, « Exemplum in Quadrato », « Exemplum in Cubo »,
%              negated affections direct and inverse, and a sentence on the
%              ancients left unfinished at the foot of view 72.
%   v. 73–74   Leaf « 34 » begins with a sentence that does not visibly
%              continue view 72 (see below); then « De reductione quadratorum
%              adfectorum ad pura. formulæ hæ. » (three formulas, numbered
%              I–III above the line) and « De reductione cuborum simpliciter
%              adfectorum sub quadrato … formulæ tres. » (I–III).
%   v. 75–78   « De reductione cuborum adfectorum tam sub quadrato quam
%              latere ad cubos adfectos simpliciter sub latere. formulæ
%              septem. » (I–VII), each with numerical examples; then, on
%              view 78, the expurgation by triangular and pyramidal unciæ.
%   v. 79–80   Leaf « 37 »: triangular and pyramidal numbers; then « De
%              transmutatione πρῶτον ἔσχατον quæ remedium est adversus vitium
%              negationis. Cap. II », with its worked examples. View 80 stops
%              mid-sentence (« itaque ex iis quæ proponuntur ») and runs on
%              into view 81.
%
% Where the batch begins and ends. View 61 continues Prop. X, whose
% enunciation ends view 60 (batch 3). View 80 runs on into view 81. Inside
% the batch one seam is unsure: view 72 ends « Credebant autem antiqui
% quoniam hac reductione æquationes quadraticas omnino expurgabant et facile
% explicabant » without a stop, and view 73 (leaf « 34 ») opens on a partly
% effaced line that the pass does not read as its continuation; a leaf may
% be missing or the sentence may run on — left open.
%
% Notation, kept as written. The unknown is A, the second root E; the given
% magnitudes carry their dimension — « B plano », « Z solido », « Z
% planosolido », « Z solidosolido », « B planoplano » — and the powers are in
% words: quadratum, cubus, quadratoquadratum, quadratocubus, cubocubus. « in »
% is multiplication; « bis », « ter » are coefficients. No sign of equality:
% « æquetur », « æquabitur », « æquatur ». Numerical examples in the cossic
% abbreviations N, Q, C, QQ, QC, CC with the coefficient before (« 1CC + 21N
% æquabitur 22 »); « est 1N » or « fit 1N » gives the root. « L », « LC »,
% « LQQ », « LQC » before a number mark the square, cube, fourth and fifth
% root, and are set $L$, $LC$ … as written. The double stroke « = » stands
% where a sign is left unspecified (a difference in either order) and is set
% « = »; where the writer crossed it into a « + » it is set \overset{+}{=}.
% Braces joining lines of a polynomial are set as arrays. Ranks I, II, III…
% written over the terms of a series are set with \overset. The marginal
% notes use their own shortened species — « Z pl. », « Z ss », « Bq »,
% « Aqq », « Ec » — and a sign of equality like a looped z, set \propto. The
% Greek words of the Emendatio are transcribed letter by letter as the hand
% writes them, without restoring accents it does not show. Underlining in the
% text is set \emph{} (text-mode \underline fails the TEI check); in
% mathematics \underline stays. The long s is set s.
%
% Corrections. The writer checks his examples: wrong numbers are underlined
% in the line and the right ones written in the margin (views 65–67, 72,
% 74–78, 80). A few errors are left uncorrected on the leaf and noted: « 30Q »
% for 30N at view 78; the equation « 1C − 12N æquatur 1684 » at view 74 is
% corrected in the margin, but « 1C + 6Q æquatur 1600 » there reads « 600 »
% with a poorly separated initial 1. Every computation cited in a \note{} was
% checked by the pass and is the pass's own.
%
% Numbers on the leaves. An ink number at the top right of each recto, one
% per leaf, versos carrying none: 28 (v. 61), 29 (63), 30 (65), 31 (67), 32
% (69), 33 (71), 34 (73), 35 (75), 36 (77), 37 (79). It continues batch 3's
% series (18 … 27 on views 41–59) without a gap, runs on across the change of
% piece at view 69, and numbers small and large leaves alike. It is in ink,
% not a pale pencil, and batch 3 found the writer referring to leaf 18 by it:
% it is taken as the writer's or a contemporary's foliation, not the
% library's, and no \folio{} is written. Figures seen in advance on
% projecting leaves (« 29 » at v. 61, « 33 » at v. 67 and 69, « 37 » at v. 73,
% 75, 77) and faint mirror-images of the recto number on versos (« 8 » at
% v. 62) are the same series. Other numbers are content: chapter numbers
% (XVIII→XIX, XIX→XX, XXI; Cap. I, Cap. II of the Emendatio), proposition and
% formula numbers, ranks over series, and « 11 » (II) and « 111 » (III)
% heading formulas. The margin's « p. 127. editionis 1646. » (v. 69) is a
% page reference to a printed edition, in a later hand.

% Coordinator's reconciliation (2026-09-23), after all nine batches:
%   One ink series of leaf numbers runs top right of the rectos through the
%   whole volume, with no gap: 1–8 and « 5 bis » (v. 5–20), 9–17 (v. 22–39),
%   18–27 (v. 41–59), 28–37 (v. 61–79), 38–47 and « 46 bis » (v. 81–100),
%   « 47 bis » and 48–56 (v. 102–120), 57–66 (v. 122–140), 67–78
%   (v. 142–160), 79 (v. 162). It crosses the three pieces of the volume, and
%   the writer cites it himself (« fol. 18. », v. 43; « pag. 48 », v. 100), so
%   it is the writer's or copyist's count of leaves, not the library's. Being
%   ink, it gets no \folio{}, in any batch.
%   The pieces: « Francisci Vietæ De Recognitione Æquationum tractatus »
%   (heading on v. 5), Cap. 1–XXI, v. 5–68; « Francisci Vietæ. De æquationum
%   Emendatione tractatus secundus » (heading on v. 69), Cap. I–XIIII, ending
%   « Finis » on v. 141; then astronomy without a name, v. 142–163 (lunar
%   prostaphaereses, Tycho's lunar hypothesis restated, « Harmonici
%   Ptolemaici Constructio Caput XII », Mercury). Two doubtful notes may mark
%   the treatises as a copy: « transcripta ex exemplari » (v. 100) and
%   « deerant hæc in exemplari » (v. 114). The name on the leaves is in the
%   headings; who wrote the hands is not settled.

\documentclass[11pt,a4paper]{article}
% The preamble shared by every transcription of the mathematicians' manuscripts in Gallica.
%
% It rests on one idea: **uncertainty compounds, it does not resolve.** A
% transcription that smooths over a doubtful word has destroyed the only thing
% it was made for — a reader must be able to tell, at every point, what was on
% the page and what was supplied. The apparatus macros below are therefore the
% critical apparatus, not decoration.
%
% The same commands are recognised by `scripts/render.mjs`, which produces the
% site's reading view, and by `scripts/tei.mjs`. A macro added here must be
% added there too, or rendering fails loudly — which is the wanted behaviour.
%
% Comments here are English, like the rest of the repository. The documents
% this preamble sets are French, and that is deliberate: see the skills.

\ProvidesPackage{germain}[2026/09/15 Transcriptions of mathematicians' manuscripts in Gallica]

\usepackage{amsmath,amssymb,amsthm}
\usepackage{mathrsfs}
% Kept from the parent preamble so the renderer's diagram support stays
% available; these manuscripts draw no commutative diagrams, and nothing here uses it.
\usepackage{tikz-cd}
\usepackage[english,french]{babel}
\usepackage{fontspec}

% --- Greek ------------------------------------------------------------------
% The text face stays fontspec's default, Latin Modern, which has no Greek:
% XeTeX drops every glyph it lacks with a line in the log and nothing on the
% page, and the Greek words the manuscripts quote vanished from the PDFs. So
% the Greek and Greek Extended blocks get a XeTeX character class of their own,
% and entering or leaving a run of it switches the family to CMU Serif — the
% same Computer Modern design, with polytonic Greek — and back. Only the
% family changes: a Greek word in \textit or \textbf keeps its shape and series.
% The transitions to and from every other class carry the tokens that class
% has towards an ordinary letter, so French spacing before « ; » or inside
% « » still holds after a Greek word. No grouping, so no brace or \endgroup
% can come out unbalanced; maths is left alone, since XeTeX inserts nothing
% there. Kept identical in `exercices.sty`.
\makeatletter
\newfontfamily\germainGreekfont{cmun}[
  Extension=.otf, NFSSFamily=germaingreek,
  UprightFont=*rm, ItalicFont=*ti, SlantedFont=*sl,
  BoldFont=*bx, BoldItalicFont=*bi, BoldSlantedFont=*bl]
\newXeTeXintercharclass\germain@greek
\def\germain@greekrange#1#2{%
  \@tempcnta=#1\relax
  \loop
    \XeTeXcharclass\@tempcnta=\germain@greek
    \ifnum\@tempcnta<#2\advance\@tempcnta\@ne\repeat}
\germain@greekrange{"0370}{"03FF}
\germain@greekrange{"1F00}{"1FFF}
\def\germain@greekprev{\rmdefault}
\def\germain@greekin{%
  \edef\germain@greekprev{\f@family}\fontfamily{germaingreek}\selectfont}
\def\germain@greekout{\fontfamily{\germain@greekprev}\selectfont}
\def\germain@greekjoin{%
  \XeTeXinterchartokenstate=\@ne
  \@tempcnta=\z@
  \loop
    \ifnum\@tempcnta=\germain@greek\else
      \XeTeXinterchartoks\germain@greek\@tempcnta=\expandafter{%
        \expandafter\germain@greekout\the\XeTeXinterchartoks\z@\@tempcnta}%
      \XeTeXinterchartoks\@tempcnta\germain@greek=\expandafter{%
        \the\XeTeXinterchartoks\@tempcnta\z@\germain@greekin}%
    \fi
    \ifnum\@tempcnta<4095\advance\@tempcnta\@ne\repeat}
% babel-french sets its own classes, and saves and restores the interchar
% state, each time a language is selected: join again after it does.
\addto\extrasfrench{\germain@greekjoin}
\addto\extrasenglish{\germain@greekjoin}
\AtBeginDocument{\germain@greekjoin}
\makeatother

\usepackage{geometry}
\geometry{a4paper,margin=25mm,marginparwidth=28mm}
\usepackage{marginnote}
\usepackage{xcolor}
\usepackage[normalem]{ulem}
\setlength{\emergencystretch}{3em}
\usepackage{hyperref}
\hypersetup{colorlinks=true,linkcolor=black,urlcolor=black,pdfborder={0 0 0}}

% --- Batch metadata ---------------------------------------------------------
% Taken as they stand by the rendering script, which shows them at the head of
% the reading view. They fix what the file claims to cover. The macro names are
% the parent project's, so that the scripts stay shared; \folder here means the
% volume's slug — `fr-9115` — and \pages counts Gallica views.

\newcommand{\folder}[1]{\gdef\grothFolder{#1}}
\newcommand{\batch}[1]{\gdef\grothBatch{#1}}
\newcommand{\pages}[2]{\gdef\grothFirst{#1}\gdef\grothLast{#2}}
\newcommand{\foldertitle}[1]{\gdef\grothFoldertitle{#1}}

% The holder's own shelfmark, printed as they write it: « Français 9115 ».
\newcommand{\shelfmark}[1]{\gdef\grothShelfmark{#1}}

% The dating, copied verbatim from the holder's catalogue — never inferred
% here. « XIXe siècle » is the BnF's; a pass that dates a leaf from its content
% says so in a \note{}, as its own inference.
\newcommand{\dating}[1]{\gdef\grothDating{#1}}

% The status notice, printed in the title block and repeated as a diagonal
% watermark on every page of the PDF. The manuscripts are in the public
% domain; what this line declares is not a right but a quality — an
% unverified first pass by a machine. The skills write the line; a file
% without it gets no watermark, so leaving it out is visible in the output.
\newcommand{\watermark}[1]{\gdef\grothWatermark{#1}}

\gdef\grothFolder{?}\gdef\grothBatch{}\gdef\grothFirst{?}\gdef\grothLast{?}%
\gdef\grothFoldertitle{}\gdef\grothShelfmark{}\gdef\grothDating{}\gdef\grothWatermark{}

% --- The résumé -------------------------------------------------------------
\newenvironment{resume}
  {\small\setlength{\parskip}{0.45em}}
  {\par\medskip\hrule\medskip}

% --- Keywords ---------------------------------------------------------------
% In English, closing the modernised reading's résumé: the modern vocabulary
% under which to search for what the leaves construct. The manifest extracts
% them to tag the volume — this line is the single source of the tags.
\newcommand{\keywords}[1]{\par\smallskip\noindent
  {\footnotesize\textsc{Keywords} --- \textit{#1}}\par}

% --- The critical apparatus -------------------------------------------------

% The Gallica view number, in the margin. This is the anchor that lets a
% disagreement over a formula be settled by looking at *one* image rather than
% a volume of 750. The site uses it to turn the facsimile as the transcription
% is scrolled. A view is one scanned image: usually one side of a leaf.
\newcommand{\page}[1]{%
  \marginnote{\footnotesize\color{gray!70}\textsf{v.\,#1}}%
  \label{page:#1}\ignorespaces}

% The BnF's pencilled foliation, where the leaf carries it and a pass can read
% it: « 198r ». This is what the literature cites — « ff. 198r–208v » — and
% the one line that turns a citation into a view. Recorded, never inferred:
% a folio a pass cannot read is not written.
\newcommand{\folio}[1]{%
  \marginnote{\footnotesize\color{gray!50}\textsf{f.\,#1}}\ignorespaces}

% The modernised reading's coarser counterpart to \page{}: which views a
% section is drawn from.
\newcommand{\pagerange}[2]{%
  \marginnote{\footnotesize\color{gray!70}\textsf{v.\,#1--#2}}%
  \ignorespaces}

% Illegible: never guessed.
\newcommand{\ill}{\textcolor{red!60!black}{[\,\ldots\,]}}

% A doubtful reading: the word is given, and flagged as uncertain.
\newcommand{\uncertain}[1]{\uline{#1}}

% An editorial addition — a missing letter, an implied word.
\newcommand{\add}[1]{\textcolor{blue!50!black}{[#1]}}

% Struck out by the writer. What was crossed out is part of the document.
\newcommand{\struck}[1]{\sout{#1}}

% The transcriber's note, on the state of the page or the mathematics.
\newcommand{\note}[1]{\par\smallskip{\small\color{gray!60!black}\textit{[#1]}}\par\smallskip}

% A marginal note of hers, distinct from the body of the text.
\newcommand{\marginal}[1]{\par\smallskip\noindent\hspace{1em}%
  {\small\color{gray!60!black}#1}\par\smallskip}

% --- Title ------------------------------------------------------------------
% The title block is French, like the documents themselves.

\AddToHook{shipout/background}{%
  \ifx\grothWatermark\empty\else
    \begin{tikzpicture}[remember picture, overlay]
      \node[rotate=52, scale=3.4, align=center, text=gray!18,
            font=\sffamily\bfseries]
        at (current page.center) {\grothWatermark};
    \end{tikzpicture}%
  \fi}

\AtBeginDocument{%
  \begin{center}
    {\large\bfseries Manuscrits de mathématiciens (Gallica) ---
      \ifx\grothShelfmark\empty\grothFolder\else\grothShelfmark\fi}\\[2pt]
    {\itshape\grothFoldertitle}\\[4pt]
    {\small vues \grothFirst--\grothLast
      \ifx\grothBatch\empty\else\ (lot \grothBatch)\fi}\\[2pt]
    \ifx\grothDating\empty\else
      {\small\color{gray!60!black}Datation du catalogue : \grothDating}\\[2pt]
    \fi
    \ifx\grothWatermark\empty\else
      {\small\bfseries\color{red!45!gray}\grothWatermark}\\[2pt]
    \fi
    {\footnotesize\color{gray!60!black}Transcription établie d'après les images IIIF de Gallica.
     Source gallica.bnf.fr / Bibliothèque nationale de France.\\
     Les lectures incertaines sont soulignées, les illisibles notées [\,\ldots\,],
     les ajouts entre crochets ; \textsf{v.} numérote les vues de Gallica,
     \textsf{f.} la foliotation au crayon de la BnF.}
  \end{center}
  \vspace{1em}}

\endinput


\folder{nal-1644}
\shelfmark{NAL 1644}
\batch{4}
\pages{61}{80}
\dating{}
\watermark{Lecture automatique\\non vérifiée}
\foldertitle{Viète (François). Papiers divers. NAL 1644}

\begin{document}

\page{61}
\note{Un petit feuillet, numéroté « 28 » à l'encre en haut à droite, posé sur un feuillet plus grand numéroté « 29 », dont on voit le haut au-dessus de lui : « prop. XV. Si $B$ plano in $A$ cubum $-\,A$ quadratocubo æquetur $Z$ planosolido. » — ces lignes sont celles de la vue 63 et y sont transcrites. La page continue la prop.~X, dont l'énoncé termine la vue précédente. Dans la marge de gauche, des traces pâles d'écriture, décharge du feuillet voisin, non transcrites.}

Est $B$ plano compositum ex duobus quadratis duorum laterum, primi ductu in secundum fit $Z$ planoplano et fit $A$ quadratum maius duorum minusve.

Sint latera 1. 4. dicetur $17Q - 1QQ$ æquari 16. et fit $1N$. 1 vel 4. quod si $1N$ intelligatur quadratum radixve plana $17N - 1Q$ æquabitur 16. et fit $1N$. 1 vel 16.
\marginal{\uncertain{$Q\,N$.} — en regard de la dernière ligne, le $N$ barré dessus et dessous.}

\subsection*{prop. XI.}

\[ \left. \begin{array}{l} \text{Si } B \text{ solidum in } A \text{ cubum} \\ -\,A \text{ cubocubo} \end{array} \right\} \text{æquetur } Z \text{ solidosolido.} \]

Est $B$ solidum compositum ex duobus cubis quorum primi ductu in secundum fit $Z$ solidosolidum. et fit $A$ cubus maior duorum minorve.

Sunto latera 1. 8. dicetur $513C - 1CC$ æquari 512 et fit $1N$. 1 vel 8.

quod si $1N$ intelligatur cubus radixve solida $513N - 1Q$ æquabitur 512 et fit $1N$ 1. vel 512.

\subsection*{prop. XII.}

\[ \left. \begin{array}{l} \text{Si } B \text{ planoplano in } A \text{ quadratum} \\ -\,A \text{ cubocubo} \end{array} \right\} \text{æquetur } Z \text{ solidosolido.} \]

Est $B$ plano compositum ex quadratis trium planorum proportionalium. et $Z$ solidosolidum quod fit ab uno plano in aggregatum quadratorum a duobus reliquis planis et fit $A$ quadratum maius minusve.
\marginal{× id est $A$q. est æquale \uncertain{quadrato} maiori minorive extremæ}

Sunto tria plana proportionalia 1. 2. 4. dicetur $21Q - 1CC$. æquari 20. et fit $1N$. 1. vel 2. quod si $1N$ intelligatur quadratum radixve plana $21N - 1C$ æquabitur 20 et fiet $1N$. 1 vel 4.
\note{La note marginale, d'une main plus petite, est appelée par une croix ; « \uncertain{quadrato} » est à demi couvert d'une tache.}

\subsection*{prop. XIII}

\[ \left. \begin{array}{l} \text{Si } B \text{ plano in } A \text{ quadratoquadratum} \\ -\,A \text{ cubocubo} \end{array} \right\} \text{æquetur } Z \text{ solidosolido} \]

Est $B$ plano compositum a tribus planis proportionalibus et $Z$ solidosolidum quod fit ab uno plano in quadratum compositi ex duobus reliquis et × fit $A$ quadratum, maius duorum minusve.
\marginal{× id est $A$q. æquale composito ex prima et secunda, vel ex secunda et tertia}

Sint plana proportionalia 5. 20. 80 dicetur $105QQ - 1CC$ æquari 50000 et fit $1N$. \emph{25}. vel \emph{100}.
\marginal{\emph{5} \quad \emph{10}}
\note{« 25 » et « 100 » sont soulignés dans la ligne, et sous eux, dans la marge de gauche, « 5 » et « 10 », soulignés aussi : ce sont les valeurs de $1N$, dont 25 et 100 sont les carrés ($105 \cdot 5^4 - 5^6 = 50000$, $105 \cdot 10^4 - 10^6 = 50000$). Le soulignement est rendu par l'italique.}

\page{62}
\note{Verso du petit feuillet « 28 » ; le texte continue celui de la vue 61 (la prop.~XIII). En haut à gauche, un « 8 » pâle, sans doute le « 28 » du recto vu par transparence. La colonne de gauche porte un calcul de la main plus petite des notes marginales, séparé du texte par un trait vertical.}

Inventio autem trium planorum proportionalium quorum medio adscito sive primo sive postremo fit quadratum numero patet ex hac serie
\[ \frac{B \text{ quadratoquadrato}}{B \text{ quad.} = G \text{ quad.}} \qquad \frac{B \text{ quad. in } G \text{ quad.}}{B \text{ quad.} = G \text{ quad.}} \qquad \frac{G \text{ quad. quad.}}{B \text{ quad.} = G \text{ quad.}} \]
\note{Le signe « $=$ » des trois dénominateurs n'est pas l'égalité : c'est sans doute le signe de la différence de deux grandeurs, d'ordre non spécifié. Dans le premier, un trait vertical traverse les deux barres, comme un « $+$ » écrit par-dessus ; dans la marge de gauche, à la même hauteur, trois « $+$ » : « $+ \quad + \quad +$ ». Le calcul qui suit (« $B$ quad. $+$ $G$ quad. ») demande en effet la somme.}

\ill{} videlicet \uncertain{totali}.

Sit enim medium planorum $E$ plano et primo statuatur $B$ quad. minus $E$ plano postremum $G$ quad. minus $E$ plano cum igitur primo plano additur $E$ planum fiet quadratum nempe $B$ quad. restat igitur ut ea tria plana proportionalia sint et consequenter quod a medio fit in se æquetur ei quod fit ab extremis qua comparatione secundum artem facta.
\[ \frac{B \text{ quad. in } G \text{ quad.}}{B \text{ quad.} + G \text{ quad.}} \quad \text{invenitur æquari } E \text{ plano} \]

\marginal{\uncertain{Duc} \ill{} $G\,p$ in $D\,p$. $-\,G\,p$ in $E\,p$. $-\,D\,p$. in $E\,p$. $+\,E\,pp$ $\mid$ $E\,pp$ — igitur $G\,p$ in $D\,p$. æquabitur $G\,p$ in $E\,p$. $+\,D\,p$ in $E\,p$. $\mid$ \uncertain{quando} — \uncertain{quibus demptis} æquantur $-E\,pp$. $\mid$ $E\,pp$. — igitur $\frac{G\,p \text{ in } D\,p.}{G\,p + D\,p}$ $\mid$ $E\,p$.}
\note{Calcul de la marge de gauche, sur deux colonnes séparées par un trait vertical : le produit $(G - E)(D - E)$ égalé à $E^2$, d'où $E = \frac{GD}{G+D}$ ; « $p$ » et « $pp$ » abrègent « plano » et « planoplano ». Les premiers mots sont mal formés ; la colonne de droite n'est lue qu'en partie.}

Sit $B$ 1. $G$ 2. plana quæsita erunt $\frac{1}{5}$. $\frac{4}{5}$. \uncertain{$\frac{16}{15}$}
\marginal{\emph{5}}
\note{Dans le troisième dénominateur, le « 1 » est un trait épais, peut-être une rature ; le calcul donne $\frac{16}{5}$, et la multiplication par 25 qui suit donne 80.}

medium adscito primo facit 1. adscito vero secundo facit 4 eadem plana ducantur in aliquod quadratum utpote 25 fient
\[ \overset{\text{I}}{5}. \quad \overset{\text{II}}{20}. \quad \overset{\text{III}}{80} \]
qualia ad exemplum adsumpta sunt.

\section*{De reliquis.}

\subsection*{Prop. XIIII}

\[ \left. \begin{array}{l} \text{Si } B \text{ solidum in } A \text{ quadratum} \\ -\,A \text{ quadratocubo} \end{array} \right\} \text{æquetur } Z \text{ planosolido.} \]

Est $B$ solidum constans cubo compositæ a duabus primis in serie trium proportionalium plus cubo compositæ a duabus postremis et insuper solido quod fit ductu alterius extremarum in quadratum compositæ ex secunda et tertia.
\marginal{1\textsuperscript{o} et \struck{\ill{}} $Z$ planosolidum quod fit \struck{\ill{}} solido \ill{} \add{\uncertain{primi}} \uncertain{cubo} a duabus \ill{} in quadratum compositæ ex prima et secunda \struck{\ill{}}}
\note{La note marginale, de la main plus petite, est appelée par le même « 1\textsuperscript{o} » que celui écrit au-dessus de « vel » à la ligne suivante ; elle fournit la définition de $Z$, que le texte ne donne pas. Elle est raturée en plusieurs endroits et tachée à la fin ; lecture très incomplète.}

\textsuperscript{1o} vel a cubo compositæ ex secunda et tertia plus solido sub tertia et quadrato compositæ ex prima et secunda\textsuperscript{a} \struck{\ill{}} \emph{in quadratum} compositæ ex prima et secunda.
\note{Après « secunda », un petit « a » suscrit, puis un signe noirci en forme de croix ; « in quadratum » est souligné. La phrase ne donne pas $B$ telle qu'elle est écrite : dans l'exemple, $B = 279 = 216 + 63$ et $63 = (4 + 3) \cdot 9$.}

et fit $A$ composita ex duabus primis vel composita ex duabus postremis.

Sint proportionales $\overset{\text{I}}{1}$. $\overset{\text{II}}{2}$. $\overset{\text{III}}{4}$. $279Q - 1QC$ æquabitur 2268. et fiet $1N$: 3 vel 6.

\page{63}
\note{Le grand feuillet numéroté « 29 » à l'encre en haut à droite, dont la vue 61 montrait le haut au-dessus du petit feuillet « 28 ». La page s'ouvre sur une proposition nouvelle ; la prop.~XIIII de la vue 62 finit avec son exemple. Dans la marge de gauche, des traces pâles, décharge ou transparence, non transcrites. Le bas du feuillet est blanc.}

\subsection*{prop. XV.}

\[ \left. \begin{array}{l} \text{Si } B \text{ plano in } A \text{ cubum} \\ -\,A \text{ quadratocubo.} \end{array} \right\} \text{æquetur } Z \text{ planosolido.} \]

Est $B$ plano constans quadrato compositæ a tribus primis in serie quatuor continue proportionalium plus quadrato compositæ e tribus postremis minus plano quod fit a tertia in compositam ex tertia secunda et prima vel secunda in compositam ex secunda tertia et quarta.

et $Z$ planosolidum quod fit a $B$ plano minus quadrato compositæ ex tribus primis in cubum compositæ \uncertain{ex} tribus \struck{postremis}. vel ab $B$ plano minus quadrato compositæ e tribus postremis in cubum compositæ ex tribus \struck{primis}
\marginal{• primis \quad postremis}
\note{Les deux mots barrés dans le texte sont corrigés dans la marge de gauche, l'un sous l'autre : « primis », « postremis ». Le calcul de l'exemple confirme la correction : $(217 - 49) \cdot 7^3 = (217 - 196) \cdot 14^3 = 57624$. Devant le premier « tribus », une tache couvre la préposition.}

et fit $A$ composita ex tribus primis vel composita ex tribus postremis.

Sint proportionales continue $\overset{\text{I}}{1}$. $\overset{\text{II}}{2}$. $\overset{\text{III}}{4}$. $\overset{\text{IIII}}{8}$.

$217C - 1QC$ æquabitur 57624 et fit $1N$. 7 vel 14.

\section*{Æqualitatum contradicentium constitutiva Cap. XVIII XIX}
\note{« XVIII », dont le V est repassé d'une encre plus noire, est suivi de « XIX » d'une encre plus pâle, sans rature de l'un ni de l'autre : un renumérotage du chapitre.}

Contradicentium autem constitutio est huiusmodi.
\note{Le trait qui traverse « est huius- » est la barre du t de « est » prolongée, non une rature.}

\subsection*{prop. I.}

\[ \left. \begin{array}{l} \text{Si } A \text{ quadratum} \\ +\,B \text{ in } A \end{array} \right\} \text{æquetur } Z \text{ plano} \]

\[ \left. \begin{array}{l} \text{et rursus } E \text{ quadratum} \\ \qquad -\,B \text{ in } E \end{array} \right\} \text{æquetur } Z \text{ plano.} \]
\note{Au-dessus de « plano », dans la seconde équation, quelques signes petits et pâles, peut-être « $+B$ in $A$ » : \ill{}.}

Sunt duo latera quorum differentia est $B$ quod autem sub eis fit æquale est $Z$ plano. et fit $A$ minus laterum \uncertain{$E$ maius}.
\marginal{\ill{}}
\note{« est $B$ » est traversé par la barre prolongée d'une lettre, comme plus haut. Les derniers mots de la ligne sont serrés contre le bord du feuillet. Dans la marge de gauche, à cette hauteur, un signe pâle suivi d'une barre oblique.}

Sunto latera 1. 2. $1Q + 1N$ æquabitur 2 et fit $1N$. 1.

rursus $1Q - 1N$. æquabitur 2. et fit $1N$. 2.

\subsection*{prop. II}

\[ \left. \begin{array}{l} \text{Si } A \text{ quadratoquadratum} \\ +\,B \text{ solido in } A \end{array} \right\} \text{æquetur } Z \text{ planoplano.} \]

\[ \left. \begin{array}{l} \text{et rursus } E \text{ quadratoquadrato} \\ \qquad -\,B \text{ solido in } E \end{array} \right\} \text{æquetur } Z \text{ planoplano} \]

\page{64}
\note{Verso du grand feuillet « 29 » ; le texte continue la prop.~II commencée au bas de la vue 63. Le petit feuillet « 28 » dépasse sous le bas de la page : on y lit la dernière ligne de la vue 62, « 2268. et fiet $1N$: 3 vel 6. », non transcrite ici. À droite, la marge du feuillet suivant, avec des notes de la main des annotations, qui reparaissent à la vue 65. En haut à gauche, une tache et un signe \ill{}. Toute la marge de gauche est occupée par des notes de la main plus petite, en trois blocs séparés par des traits.}

Sunt quatuor continue proportionales a quibus cubi alterne sumpti differunt per $B$ solidum. fit autem $Z$ planoplano ductu alterius extremæ in differentiam cuborum a reliquis alterne sumptorum et est $A$ prima minor inter extremas $E$ quarta.

Sunto continue proportionales. $\overset{\text{I}}{1}$. $\overset{\text{II}}{LC2}$. $\overset{\text{III}}{LC4}$. $\overset{\text{IIII}}{2}$.
\note{« $LC$ » : \emph{latus cubicum}, la racine cubique ; de même « $LQC$ » plus bas, \emph{latus quadratocubicum}, la racine cinquième.}

$1QQ + 5N$. æquabitur 6. et fit $1N$. 1.

rursus $1QQ - 5N$. æquabitur 6. et fit $1N$. 2.

\marginal{4 Sit differentia alterne sumpta trium primarum $A$. ergo $\frac{Z\,pp}{A}$ erit \uncertain{quarta} maior æquale $B + A$. sive differentiæ quatuor continue proportionalium alterne sumptarum plus differentia trium primarum alterne sumptarum ut ex hac serie patet. \quad $A$. $B$. $\frac{Bq}{A}$ $\frac{Bc}{Aq}$. \quad differentia quatuor alterne sumptarum est $\frac{Bc \text{ in } A + B \text{ in } Ac - Bq \text{ in } Aq - Aqq}{Ac}$ \quad differentia trium primarum sumptarum alterne est $\frac{Bq + Aq - B \text{ in } A}{A}$. \quad quare $\frac{Bc \text{ in } Aq + B \text{ in } Aqq - Bq \text{ in } Ac - Aqc + Bq \text{ in } Ac + Aqc - B \text{ in } Aqq}{Aqq}$ æquale erit $\frac{Bc}{Aq}$}
\note{Premier bloc de la marge de gauche, en regard de la fin de la prop.~II et de la prop.~III : une démonstration par la suite $A$, $B$, $\frac{Bq}{A}$, $\frac{Bc}{Aq}$ ; « $q$ », « $c$ », « $qq$ », « $qc$ » abrègent quadratum, cubus, etc. Les deux sommes sont écrites en colonne, terme sous terme. Dans le dernier terme du numérateur, « $Aqq$ » est surchargé et peu sûr. Le « 4 » en tête est le numéro de la remarque.}

\subsection*{prop. III}

\[ \left. \begin{array}{l} \text{Si } A \text{ cubocubus} \\ +\,B \text{ planosolido in } A \end{array} \right\} \text{æquetur } Z \text{ solidosolido} \]

\[ \left. \begin{array}{l} \text{et rursus } E \text{ cubocubus} \\ \qquad -\,B \text{ planosolido in } E \end{array} \right\} \text{æquetur } Z \text{ solidosolido.} \]

Sunt sex continue proportionales a quibus quadratocubi alterne sumpti differunt per $B$ planosolidum. fit autem $Z$ solidosolidum ductu alterius extremæ in differentiam quadratocuborum a reliquis alterne sumptorum et est $A$ prima minor inter extremas. $E$ sexta.

Sunto proportionales continue
\[ \overset{\text{I}}{1}.\ \overset{\text{II}}{LQC2}.\ \overset{\text{III}}{LQC4}.\ \overset{\text{IIII}}{LQC8}.\ \overset{\text{V}}{LQC16}.\ \overset{\text{VI}}{2}. \]

$1CC + 21N$. æquabitur 22. et fit $1N$. 1.

et rursus $1CC - 21N$. æquabitur 22. et fit $1N$. 2.

\subsection*{prop. IIII.}

\[ \left. \begin{array}{l} \text{Si } A \text{ quadratoquadratum} \\ +\,B \text{ in } A \text{ cubum} \end{array} \right\} \text{æquetur } Z \text{ planoplano} \]

\[ \left. \begin{array}{l} \text{et rursus } E \text{ quadratoquadratum} \\ \qquad -\,B \text{ in } E \text{ cubum} \end{array} \right\} \text{æquetur } Z \text{ planoplano,} \]

Sunt quatuor continue proportionales quarum alterne sumptarum × dum intelligitur \struck{$B$} prima, minor inter extremas (et) sit $A$ differentia alterne sumpta trium primarum $E$ differentia trium postremarum.

\marginal{Eodem \add{modo} ostenderetur si \uncertain{fingerentur} \ill{} quatuor continue proportionalium cum intelligitur prima minor, differentiam trium postremarum alterne sumptarum \uncertain{minus} differentia quatuor alterne sumptarum \emph{æquari extremæ \uncertain{maiori}}}
\note{Deuxième bloc de la marge, encadré de traits au-dessus et au-dessous ; il se rapporte à la seconde équation de la prop.~II ou de la prop.~IIII, sans renvoi visible. « modo » est écrit au-dessus de la ligne ; la dernière ligne est soulignée.}

\marginal{× differentia est $B$ et fit $Z$ planoplanum ex ductu \uncertain{utriusque} inter extremas in cubum differentiæ reliquarum alterne sumptarum. et \ill{} est namque $\frac{1}{4} \big/ 5$ et rursus $\frac{2}{8} \big/ 10$ quarum differentia est 5.}
\note{Troisième bloc de la marge, appelé par la croix écrite dans la ligne après « sumptarum » : il complète la phrase, que le texte laisse sans verbe. Les nombres de la fin sont écrits en colonne : 1 sur 4, et leur somme 5 sous un trait ; 2 sur 8, et 10. Dans l'exemple, $1 + 4 = 5$ et $2 + 8 = 10$ diffèrent de $B = 5$, et $Z = 216 = 1 \cdot 6^3 = 8 \cdot 3^3$.}

Sunto proportionales continue $\overset{\text{I}}{1}$. $\overset{\text{II}}{2}$. $\overset{\text{III}}{4}$. $\overset{\text{IIII}}{8}$.

$1QQ + 5C$. æquabitur 216. et fiet $1N$. 3

et rursus $1QQ - 5C$. æquabitur 216 et fiet $1N$. 6.

\marginal{In secunda æquatione $\frac{Z\,pp}{Ec} \propto E - B$. id est differentia trium postremarum \struck{\ill{}} alterne sumptarum minus differentia quatuor alterne sumptarum}
\note{Cette dernière note, de la main plus petite, est écrite sous l'exemple, dans la colonne du texte, ouverte par une grande parenthèse. Le signe d'égalité de cette main, une boucle en forme de z, est rendu $\propto$ ; le « $c$ » du dénominateur est petit. Dans l'exemple, $\frac{216}{6^3} = 1 = 6 - 5$.}

\page{65}
\note{Feuillet numéroté « 30 » à l'encre en haut à droite. Le texte continue sans rupture la série de la vue 64. À gauche dépasse la marge du feuillet « 29 » (fins de lignes de la vue 64) ; à droite, le bord d'un autre feuillet. La marge de gauche est couverte de notes et de calculs de la main plus petite, en blocs séparés par des traits. Le bas du feuillet est blanc.}

\subsection*{Prop. V.}
\note{Le numéro est en partie couvert d'une tache : « Prop. V » se lit par la suite (prop.~IIII à la vue 64).}

\[ \left. \begin{array}{l} \text{Si } A \text{ cubocubus} \\ +\,B \text{ in } A \text{ quadratocubum} \end{array} \right\} \text{æquetur } Z \text{ solidosolido.} \]

\[ \left. \begin{array}{l} \text{et rursus} \\ \quad E \text{ cubocubus} \\ \quad -\,B \text{ in } E \text{ quadratocubum} \end{array} \right\} \text{æquetur } Z \text{ solidosolido} \]

Sunt sex continue proportionales quarum alterne sumptarum differentia est $B$. fit autem $Z$ solidosolidum ductu utriusque extremæ in quadratocubum differentiæ reliquarum alterne sumptarum

× \struck{et} dum intelligitur prima, minor inter extremas sit $A$ differentia alterne sumpta quinque primarum \struck{et} $E$ differentia quinque postremarum.

Sunto proportionales continue $\overset{\text{I}}{1}$. $\overset{\text{II}}{2}$. $\overset{\text{III}}{4}$. $\overset{\text{IIII}}{8}$. $\overset{\text{V}}{16}$. $\overset{\text{VI}}{32}$

$1CC + 21QC$. æquabitur \emph{4,618,812}, et fit $1N$. 11 et rursus

$1CC - 21QC$ æquabitur \emph{4,618,812}, et fit $1N$ \emph{22}.
\note{Les deux « 4,618,812 » sont soulignés et surlignés, et « 22 » souligné : le nombre est faux. La marge donne le bon, 5153632 : $11^6 + 21 \cdot 11^5 = 22^6 - 21 \cdot 22^5 = 5153632$. Les virgules séparent les tranches de trois chiffres.}

\marginal{$\frac{Z\,ss}{Aqc} \propto B + A$. in serie enim \struck{\ill{}} \add{sex} continue proportionalium differentia omnium alterne sumptarum plus differentia quinque primarum æquatur maiori extremæ \quad $\frac{Z\,ss}{Eqc} \propto E - B$. rursus differentia quinque postremarum minus differentia \struck{\ill{}} sex alterne sumptarum æquatur extremæ minori \uncertain{ut ex calculo} clarum est.}
\marginal{$\frac{1,\ 4,\ 16}{21}$ \quad $\frac{2,\ 8,\ 32}{42}$ \quad $\frac{42 - 21}{21.\ B.}$ \quad fit una extremarum. 1 \quad $\frac{+2,\ 8,\ 32}{42}$ \quad $\frac{4,\ 16}{20}$ \quad $\frac{42 - 20}{22.}$ cuius quadratocubus est. 5153632. \quad \struck{515363\ill{}} \quad \emph{5153632.} \struck{\ill{}} \quad × $\frac{1,\ 4,\ 16}{21}$ \quad $\frac{2,\ 8}{10}$ \quad $\frac{21 - 10}{11.}$ $A$) \quad $\frac{2,\ 8,\ 32}{42}$ \quad $\frac{4,\ 16}{20}$ \quad $\frac{42 - 20}{22.}$}
\note{Calculs de la marge, écrits en colonnes, chaque somme sous un trait : les termes de rang impair et de rang pair de la suite, leur différence $B = 21$, puis $A = 11$ et $E = 22$. Le bloc du haut est séparé des colonnes par un trait courbe. Sous « cuius quadratocubus est 5153632 », deux autres copies du nombre, la première barrée avec son dernier chiffre, la seconde soulignée deux fois et suivie d'une surcharge noircie. La croix de ce bloc répond à celle du texte, devant « dum intelligitur ».}

\section*{Æqualitatum Inversarum Constitutiva. Cap. \struck{XIX} XX}

Inversarum denique systatica sunt hæc.

\subsection*{prop. I}

\[ \left. \begin{array}{l} \text{Si } B \text{ plano in } A \\ -\,A \text{ cubo.} \end{array} \right\} \text{æquetur } Z \text{ solido} \]

\[ \left. \begin{array}{l} \text{et rursus } \underline{\overline{B}} \text{ cubus} \\ \qquad -\,B \text{ plano in } E \end{array} \right\} \text{æquetur } Z \text{ solido} \]
\marginal{$E$}
\note{« $B$ cubus » : le $B$ est souligné et surligné, et la marge porte, encadré, « $E$ », qui est la bonne lettre.}

Sunt tres proportionales a quibus quadrata alterne sumpta differunt per $B$ planum. fit autem $Z$ solidum ductu alterius extremæ in differentiam quadratorum a reliquis et est $A$ prima minor inter extremas $E$ tertia.

\marginal{$\frac{Z\,s}{A} \propto Bp - Aq$. \add{\uncertain{signandi}} differentia \struck{\ill{}} quadratorum alterne sumptorum auferatur quadratum primæ relinquitur differentia quadratorum a reliquis \struck{\ill{}} \add{\uncertain{postremis}} ut patet ex proposita serie $A$. $B$. $\frac{Bq}{A}$ \quad \uncertain{Ito} $\frac{Z\,s}{E} \propto Eq - Bp$. \uncertain{Item} a \struck{\ill{}} quadrato maioris extremæ si auferatur \struck{quadratum} differentia quadratorum a tribus alterne sumptorum relinquitur differentia quadratorum a duabus primis \struck{\ill{}} \uncertain{potissimis}.}
\note{Dernier bloc de la marge, en regard de la prop.~I ; la seconde moitié est enfermée dans une accolade tournée vers le texte. Plusieurs mots barrés, et des ajouts interlinéaires peu lisibles.}

Sunto proportionales $\overset{\text{I}}{1}$. $\overset{\text{II}}{L2}$. $\overset{\text{III}}{2}$.

$3N - 1C$ æquabitur 2. et fit $1N$. 1. et rursus

$1C - 3N$ æquabitur 2. et fit $1N$. 2.

\subsection*{prop. II.}

\[ \left. \begin{array}{l} \text{Si } B \text{ planoplano in } A \\ -\,A \text{ quadratocubo} \end{array} \right\} \text{æquetur } Z \text{ planosolido} \]

\[ \left. \begin{array}{l} \text{et rursus} \\ \quad E \text{ quadratocubus} \\ \quad -\,B \text{ planoplano in } A.\,E \end{array} \right\} \text{æquetur } Z \text{ planosolido.} \]
\note{« in $A$. $E$ » : le $E$ suit le $A$ sans que celui-ci soit barré ; c'est $E$ que l'équation demande.}

\page{66}
\note{Verso du feuillet « 30 » ; la prop.~II des \emph{inversæ}, dont l'énoncé termine la vue 65, se poursuit. Sous le bas de la page dépasse de nouveau le petit feuillet « 28 », avec la ligne « 2268. et fiet $1N$: 3 vel 6. » de la vue 62, non transcrite ici. En haut à gauche, une tache. La marge de gauche et le haut de la page portent des notes de la main plus petite ; celle du haut court sur toute la largeur, au-dessus de la première ligne du texte, qu'elle souligne d'un trait.}

\marginal{$\frac{Z\,pl}{A} \propto Bpp - Aqq$. (a differentia quadrato-quadratorum alterne sumptorum ablato quadquad. primæ relinquit differentiam quadratoquadratorum a quatuor postremis. et contra a quadquad extremæ maioris ablata differentia quadquad alterne sumptorum relinquitur diff. quadquad a quatuor primis alterne sumptorum \quad $\frac{Z\,pl}{E} \propto Eqq - Bpp$.}

Sunt quinque continue proportionales longitudines a quibus quadratoquadrata alterne sumpta differunt per $B$ planoplanum. fit autem $Z$ planosolidum ductu alterius extremæ in differentiam quadratoquadratorum a reliquis et fit $A$ prima minor inter extremas $E$ quinta.
\note{Un trait traverse « et fit $A$ » : rature ou barre de t prolongée, on ne peut trancher.}

\marginal{$B$ alterne sumptis \quad $\frac{1,\ 4,\ 16}{21}$ \quad $\frac{2,\ 8}{10}$ \quad $\frac{21 - 10}{11}$ \quad \ill{}}

Sunto continue proportionales $\overset{\text{I}}{1}$. $\overset{\text{II}}{LQQ2}$. $\overset{\text{III}}{LQQ4}$. $\overset{\text{IIII}}{LQQ8}$. $\overset{\text{V}}{2}$.

$11N - 1QC$ æquatur 10 et fit $1N$ 1. et rursus

$\underline{\overline{1QQ}} - 11N$ æquatur 10 et fit $1N$. 2.
\marginal{$\overline{\underline{C}}$}
\note{« $1QQ$ » est souligné et surligné, son second signe surchargé ; la marge porte « $C$ » entre deux traits : il faut $1QC$, comme à la ligne précédente ($2^5 - 22 = 10$).}

\marginal{$\frac{Z\,s}{Aq} \propto B + A$ differentia trium proportionalium alterne sumptorum plus differentia duarum primarum æquatur extremæ maiori. \quad $\frac{Z\,s}{Eq} \propto B - E$. si a differentia trium alterne sumptarum auferatur differentia duarum postremarum relinquitur extrema minor. \uncertain{ex} proposita serie $A$. $B$. $\frac{Bq}{A}$ \quad $\frac{4}{5}$ \quad $-\frac{2}{3}$}
\note{Bloc de la marge en regard de la prop.~III ; les deux fractions de la fin, sous un trait courbe, ne se rattachent à rien de lisible.}

\subsection*{prop. III.}

\[ \left. \begin{array}{l} \text{Si } B \text{ in } A \text{ quadratum} \\ +\,A \text{ cubo.} \end{array} \right\} \text{æquetur } Z \text{ solido.} \]

\[ \left. \begin{array}{l} \text{et rursus} \\ \quad B \text{ in } E \text{ quadratum} \\ \quad -\,E \text{ cubo} \end{array} \right\} \text{æquetur } Z \text{ solido} \]

Sunt tres proportionales quarum alterne sumptarum differentia est $B$. fit autem $Z$ solidum ductu alterutrius extremæ in quadratum differentiæ reliquarum alterne sumptarum et dum intelligitur prima minor inter extremas est $A$ differentia alterne sumpta duarum primarum $E$ differentia duarum postremarum.
\note{« tres » est écrit sur un mot effacé.}

\marginal{differentia 1 et 2 est 1 cuius q. est 1. quod ducto in 4 facit 4. sit differentia 2 et 4 est 2 cuius q. est 4. quod ducto in 1 facit 4.}

Sunto proportionales $\overset{\text{I}}{1}$. $\overset{\text{II}}{2}$. $\overset{\text{III}}{4}$.

$3Q \mathbin{\overset{+}{=}} 1C$. æquatur 4 et fit $1N$ 1. et rursus

$3Q - 1C$ æquatur 4. et fit $1N$. 2.
\marginal{$+$}
\note{Dans la première équation, le signe est un « $=$ » traversé d'un trait vertical, comme à la vue 62 ; la marge porte « $+$ » entre deux traits. L'équation demande $+$ : $3 + 1 = 4$. Le « 4 » qui suit « æquatur » est repassé.}

\subsection*{prop. IIII.}

\[ \left. \begin{array}{l} \text{Si } B \text{ in } A \text{ quadratoquadratum} \\ +\,A \text{ quadratocubo} \end{array} \right\} \text{æquetur } Z \text{ planosolido} \]

\[ \left. \begin{array}{l} \text{et rursus} \\ \quad B \text{ in } E \text{ quadratoquadratum} \\ \quad E \text{ quadratocubo} \end{array} \right\} \text{æquetur } Z \text{ planosolido.} \]
\note{Aucun signe n'est visible devant « $E$ quadratocubo ».}

Sunt quinque proportionales quarum alterne sumptarum differentia est $B$. fit autem $Z$ planosolidum ductu alterutrius extremæ in quadroquadratum differentiæ reliquarum alterne sumptarum.

et dum intelligitur prima minor inter extremas est $A$ differentia alterne sumpta quatuor primarum $E$ differentia quatuor postremarum.

\marginal{$\frac{Z\,pl}{Aqq} \propto B + A$. differentia quinque proportionalium continue plus \struck{differentia} alterne sumpta quatuor primarum æquatur extremæ maiori et rursus \quad $\frac{Z\,pl}{Eqq} \propto B - E$. si a differentia quinque continue proportionalium alterne sumptarum auferatur differentia quatuor postremarum alterne sumptarum relinquitur prima.}
\note{Derniers blocs de la marge, en regard de la prop.~IIII, le premier fermé par un trait courbe. La proposition n'a pas d'exemple ; le reste de la page est blanc.}

\page{67}
\note{Un feuillet numéroté « 31 » à l'encre en haut à droite, posé sur un feuillet plus grand numéroté « 33 », dont on voit les trois premières lignes au-dessus de lui : « Imminuis ab adiectione sub gradu extremo ut opus comprobabit postremo fit homogeneum sub $B$ coefficiente et $A$ extremo gradu adfectum \ill{} \ill{} $A$ \ill{} quoniam igitur » — lignes d'une autre pièce, lues ici sans assurance et non transcrites comme texte. Le feuillet « 31 » continue la prop.~IIII de la vue 66 par son exemple.}

\marginal{$\frac{1,\ 4,\ 16}{21.}$ \quad $\frac{2,\ 8}{10}$ \quad $\frac{21 - 10}{11}$ $B$ \quad $\frac{2,\ 8}{10}$ \quad $\frac{4,\ 16}{20}$ \quad $\frac{20 - 10}{10}$ cuius quadratoquadrat. 10000}
\note{Calculs de la marge : $B = 21 - 10 = 11$, puis $E = 20 - 10 = 10$ et $10^4$. Après « 10000 », quelques chiffres serrés, \ill{}.}

Sunto proportionales $\overset{\text{I}}{1}$. $\overset{\text{II}}{2}$. $\overset{\text{III}}{4}$. $\overset{\text{IIII}}{8}$. $\overset{\text{V}}{16}$.

$11QQ + 1QC$. æquatur 10,000. et fit $1N$ 5. et rursus

$11QQ - 1QC$. æquatur 10,000 et fit $1N$. 10.

\subsection*{prop. V.}

\[ \left. \begin{array}{l} \text{Si } B \text{ plano in } A \text{ cubum} \\ = A \text{ quadratocubo} \end{array} \right\} \text{æquetur } Z \text{ planosolido} \]
\marginal{$\overset{+}{=}$ \quad et rursus $B\,p$ in \uncertain{$E$} \quad $-\,E$ quad. cub. \quad æquetur $Z$ planosolido}
\note{Le signe devant « $A$ quadratocubo » est de nouveau un « $=$ » ; la marge porte au-dessus un « $+$ » sur deux barres, puis la seconde équation, que le texte omet. Le bloc est fermé par un trait.}

Sunt sex proportionales quarum extremæ ductæ in differentiam quartæ et primæ faciunt $B$ planum. fit autem $Z$ planosolidum e cubo differentiæ quartæ et primæ in $B$ planum plus eiusdem differentiæ quadrato. vel e cubo differentiæ quintæ et secundæ in $B$ planum, multato illius differentiæ inter quintam et secundam quadrato

et cum prima intelligitur \emph{maior} inter extremas sit $A$ differentia quartæ et primæ et × differentia quintæ et secundæ.
\marginal{\emph{minor} \quad × $E$.}
\note{« maior » est souligné et corrigé en « minor » dans la marge ; la croix après « et » appelle « $E$ », écrit dans la marge. Dans l'exemple, $Z = 7^3 \cdot 231 + 7^5 = 14^3 \cdot 231 - 14^5 = 96040$ : les deux « quadrato » de la phrase précédente valent « quadratocubo ».}

Sunto proportionales. $\overset{\text{I}}{1}$. $\overset{\text{II}}{2}$. $\overset{\text{III}}{4}$. $\overset{\text{IIII}}{8}$. $\overset{\text{V}}{16}$. $\overset{\text{VI}}{32}$.

$231\underline{Q} + 1QC$. æquat. 96040. et fit $1N$ 7. et rursus

$231\underline{Q} - 1QC$. æquat. 96040 et fit $1N$. 14.
\marginal{$\overline{\underline{C}}$ \quad $\overline{\underline{C}}$}
\note{Les deux « $Q$ » soulignés sont corrigés en « $C$ » dans la marge : $231 \cdot 7^3 + 7^5 = 96040$.}

\subsection*{prop. VI.}

\[ \left. \begin{array}{l} \text{Si } B \text{ solidum in } A \text{ quadratum.} \\ +\,A \text{ quadratocubo} \end{array} \right\} \text{æquetur } Z \text{ planosolido} \]
\marginal{× et rursus $B$ solidum in $E$ q. \quad $-\,E$ quad. cubo \quad æquetur $Z$ planosolido}
\note{La seconde équation est de nouveau suppléée dans la marge, appelée par une croix sous la première.}

Sunt sex proportionales continue quarum extremæ ductæ in quadratum a differentia tertiæ et primæ faciunt $B$ solidum. fit autem $Z$ planoplano a $B$ solido plus cubo e differentia tertiæ et primæ in quadratum differentiæ eiusdem, vel $B$ solido minus cubo e differentia secundæ et quartæ in quadratum differentiæ illius inter \emph{tertiam} et quartam.
\marginal{\emph{secundam}.}
\note{« tertiam » est souligné, corrigé en « secundam » dans la marge : $(297 - 6^3) \cdot 6^2 = 2916$.}

et cum prima intelligitur minor inter extremas sit $A$ tertia minus prima. $E$ quarta minus secunda.

Sunto proportionales. $\overset{\text{I}}{1}$ $\overset{\text{II}}{2}$. $\overset{\text{III}}{4}$. $\overset{\text{IIII}}{8}$. $\overset{\text{V}}{16}$. $\overset{\text{VI}}{32}$.

$297\underline{Q} + 1QC$. æquat. 2916. et fit $1N$. 3. et rursus

$297\underline{Q} - 1QC$. æquat. 2916. et fit $1N$. 6.
\note{Ici les « $Q$ » soulignés sont justes ($297 \cdot 9 + 243 = 2916$) et la marge ne corrige rien. Le bas de la page est blanc.}

\page{68}
\note{Verso du feuillet « 31 », sans numéro. Écriture de la même main que le recto, sans notes marginales ; à gauche, des traces pâles d'écriture, décharge du feuillet voisin, non transcrites. Le bas de la page est blanc.}

\section*{Alia rursus æqualitatum Inversarum Constitutiva Cap. XXI}

\subsection*{prop. I.}

\[ \left. \begin{array}{l} \text{Si } B \text{ plano in } A \\ -\,A \text{ cubo} \end{array} \right\} \text{æquetur } Z \text{ solido.} \]

\[ \left. \begin{array}{l} \text{et rursus} \\ \quad E \text{ cubus} \\ \quad -\,B \text{ plano in } E \end{array} \right\} \text{æquetur } Z \text{ solido.} \]

Sunt tres proportionales quarum quadrata \uncertain{iuncta} conficiunt $B$ planum. summa autem extremitatum ducta in mediæ quadratum vel altera extremarum in aggregatum quadratorum \struck{ex} duabus reliquis facit $Z$ solidum. et fit $A$ alterutra extremarum $E$ vero earundem summa.

Sunto proportionales $1$. $2$. $4$.

$21N - 1C$. æquatur 20 et fit $1N$. 1 vel 4. et rursus

$1C - 21N$. æquatur 20 et fit $1N$. 5.

\subsection*{prop. II.}

\[ \left. \begin{array}{l} \text{Si } B \text{ in } A \text{ quadratum} \\ -\,A \text{ cubo} \end{array} \right\} \text{æquetur } Z \text{ solido} \]

\[ \left. \begin{array}{l} \text{et rursus} \\ \quad B \text{ in } E \text{ quadratum} \\ \quad +\,E \text{ cubo} \end{array} \right\} \text{æquetur } Z \text{ solido} \]

Sunt tres proportionales radices quarum summa est $B$ composita autem e duabus primis \uncertain{adducta} compositæ e duabus reliquis, dum ducitur in mediæ quadratum vel altera extremarum in quadratum compositæ e duabus reliquis facit $Z$ solidum et fit $A$ alterutra prædictarum compositarum. $E$ media inter extremas.
\note{« summa » est surchargé. La première des deux constructions de $Z$ ne donne pas 36 telle qu'elle est lue ; la seconde, $1 \cdot 6^2 = 4 \cdot 3^2 = 36$, s'accorde avec l'exemple.}

Sunto proportionales $1$. $2$. $4$.

$7Q - 1C$. æquatur 36 et fit $1N$ 3 vel 6. et rursus

$7Q + 1C$. æquatur 36. et fit $1N$. 2.

\page{69}
\note{Un feuillet numéroté « 32 » à l'encre en haut à droite, posé sur le grand feuillet « 33 », dont les trois premières lignes (« Imminuis ab adiectione sub gradu extremo … quoniam igitur ») dépassent au-dessus, comme à la vue 67. Une pièce nouvelle commence ici, en tête de page, dans la même main que les feuillets précédents ; la vue 68 finissait sur l'exemple d'une proposition. À gauche, le bord d'un autre feuillet avec des fins de lignes, non transcrites.}

\section*{Francisci Vietæ. De æquationum Emendatione tractatus secundus.}

\subsection*{De solennibus quinque modis præparandarum æquationum adversus earum in numeris δυσμηχανίαν.}

\subsection*{Ac primum. De expurgatione per uncias quæ remedium est adversus πολυπλοκείαν Cap. I.}
\note{Les deux mots grecs sont écrits d'une main peu sûre, sans accents lisibles ; ils sont transcrits lettre à lettre.}

\marginal{p. 127. editionis 1646.}
\note{Dans la marge de gauche, en regard de la première ligne du chapitre, d'une autre main, petite et régulière, plus récente que le texte : « p. 127. editionis 1646. », un renvoi à une édition imprimée.}

Agnita æquationum constitutione analysta ad præparandum eas quæ sunt alioqui \uncertain{mechanicismo} respuant aut demum ægre subeunt toto se confert et qua præparatione efficit \uncertain{εὐμηχάνους}. at præparandi quidem generale doctrina est, ut nova \ill{} substituatur, vel plasmatis aut syncriseos vestigia \uncertain{repetantur}, at denique nullus non tentetur transmutandi modus. sed non desunt analystæ singularia, et topica remedia adversus vitia quæque æquationum impedimentaque quo minus facilius sive re sive numero explicentur. fere autem quæ prosunt geometricæ ad εὐμηχανίαν prosunt et arithmeticæ ut etiam e contra. et etiam de effectionibus geometricis dicetur specialius suo loco, nunc autem circa numerosam analysim magis esse intentum nostri est instituti.

Præparationum igitur præsertim in numeris solennes et speciales modi sunt fere quinque.
\begin{enumerate}
\item Expurgatio per uncias
\item transmutatio πρώτον ἔχατον
\item Anastrophe.
\item Isomæria
\item Climactica paraplerosis.
\end{enumerate}
\note{La liste est numérotée 1 à 5 par l'auteur. Au 2, les deux mots grecs sont écrits ainsi ; le second a un accent au-dessus de la première lettre.}

omnium adversus πολυπλοκείαν tutissimum et paratissimum remedium est expurgatio per uncias.

est autem species transmutationis per additionem vel subductionem.
\note{Le reste de la page est blanc ; le chapitre continue au verso, vue 70.}

\page{70}
\note{Verso du feuillet « 32 », sans numéro ; le chapitre I du \emph{De æquationum emendatione} continue. La dernière ligne s'arrête sur « quæ omnino evenit », suivi d'un mot barré : la phrase reprend au haut du feuillet « 33 » (vue 71), dont les trois premières lignes se voyaient déjà aux vues 67 et 69 au-dessus des feuillets plus petits.}

Hæ æquationes potestatum adfectarum sub extremo parodico gradu quem sustinet coefficiens radici homogenea, \ill{} ea adfectione liberantur salva numerorum symmetria. et in quadratis est \uncertain{Diahemysis}, in cubis diatritemorion, \struck{et sic in infinitum ordine} \add{in quadratoquadratis diatetartemorion.} in quadratocubis diapentemorion, in cubocubis diectemorion et sic in infinitum ordine, quoniam adfectiones sub eo gradu proficiscuntur a \ill{}, de \ill{} quo intelligit adfecta radix de qua cumque potestate pura vel affecta primo proponebatur æquatio.
\marginal{1\textsuperscript{o} in quadratoquadratis diatetartemorion.}
\note{« in quadratoquadratis diatetartemorion » est écrit dans la marge de gauche, appelé par « 1\textsuperscript{o} », et souligné d'un trait qui rejoint la ligne ; il remplace les mots barrés « et sic in infinitum ordine », que l'auteur répète plus loin à leur place. Les deux mots rendus \ill{} (« a \ill{}, de \ill{} ») sont de même racine, le second avec la terminaison -one ; aucune lecture ne s'impose.}

\uncertain{Atqui} autem \ill{} de \ill{} duplum est coefficiens sub latere in quadratis. et triplum, coefficiens sub quadrato in cubis, et quadruplum coefficiens sub cubo in quadratoquadratis. et quintuplum coefficiens \struck{coefficiens} sub quadratoquadrato in quadratocubis, et sextuplum denique coefficiens sub quadratocubo in cubocubis. et sic in infinitum ordine.

Itaque ad delendum plasma, contraria retrogradaque via fiunt unciæ conditionariæ coefficientis radici homogeneorum. in quadratis videlicet semis. in cubis triens. in quadratoquadratis quadrans. In quadratocubis quintans, in cubocubis sextans et eo continuo ordine, atque illis unciis conditionariis adficitur radix, atque adeo certe fit transmutatio.

totius itaque operis structura tres casuum differentias admittit primo.

Sit $A$ potestas adfecta per adiunctionem homogenei sub $B$ coefficiente et $A$ parodico extremo gradu, quoniam igitur ea adfectio adfirmata est, adficitur radix affirmate a conditionariis $B$ coefficientis unciis et ita affecta statuitur esse $E$, unde $E$ minus conditionariis $B$ coefficientis unciis æquabitur $A$ sub qua nova specie æquatio primo proposita \uncertain{ducetur} et ordinabitur nova quæ omnino evenit immunis ab affectione sub gradu extremo ut opus comprobabit.

Secundo sit $A$ potestas affecta per multam homogenei sub $B$ coefficiente et $A$ extremo parodico gradu, quoniam igitur ea adfectio negata est, adficitur $A$ negate a conditionariis $B$ coefficientis unciis, et ita adfecta statuitur esse $E$ unde $E$ plus conditionariis illis \struck{\ill{}} coefficientis unciis æquabitur $A$ sub qua nova specie æquatio primo proposita \uncertain{ducetur} et ordinabitur nova. quæ omnino evenit \struck{immunis}
\marginal{$B$}
\note{La lettre barrée avant « coefficientis » est corrigée dans la marge de gauche en « $B$ », entre deux traits. Le dernier mot, « immunis », est barré ; il est récrit en tête du feuillet suivant.}

\page{71}
\note{Le grand feuillet numéroté « 33 » à l'encre en haut à droite. Sa première ligne reprend le mot barré à la fin de la vue 70 et continue la phrase. À gauche dépassent les bords des feuillets « 31 » et « 32 », avec des fins de lignes, non transcrites. Sous le bas de la page dépasse un feuillet plus long, qui porte « $-\,D$ plano in $A$ » ; c'est la dernière ligne du feuillet « 35 » (vue 75), non transcrite ici.}

immunis ab adfectione sub gradu extremo ut opus comprobabit.

postremo sit homogeneum sub $B$ coefficiente et $A$ extremo gradu adfectum per multam $A$ potestatis. quoniam igitur potestas potius adficit quam adficitur utpote quæ negatur de homogeneo, adfectionis conditionariæ $B$ coefficientis unciæ multabuntur $A$ potestate et statuentur esse $E$, unde eædem coefficientis unciæ minus radice $E$, æquabuntur $A$, sub qua nova specie æquatio primo proposita \uncertain{ducetur} et ordinabitur nova quæ omnino evenit immunis ab affectione sub gradu extremo ut opus comprobabit

\subsection*{Exemplum in Quadrato.}

proponatur $A$ quadratum plus $B$ in $A$ bis æquari $Z$ plano quoniam igitur in climacticarum magnitudinum ordine latus quadrato proxime succedit, proponitur autem hic quadratum adfectum sub latere, omnino æqualitas plasmatica est alicunde effecta. est autem adfectio illa adfirmata. Itaque plasma fuit per additionem semissis coefficientis sub latere, prout conditio quadrati quæ potestas est rationis duplicatæ exposcit. ad tollendum igitur plasma, fiat expurgatio \uncertain{diahemyseos} et idcirco $A + B$ esto $E$. ergo $E - B$ \struck{\ill{}} erit $A$. et consequenter quadratum ab $E - B$ adiunctum plano sub $B$ bis et $E - B$ æquabitur $Z$ plano per ea quæ proponuntur. Æqualitas igitur de $E$ secundum artem ordinetur: omnibus rite peractis deprehendetur $E$ quadratum æquari $Z$ plano plus $B$ quadrato, quæ quidem nova æquatio pura est ab affectione sub latere, qua primo proposita obruebatur. cum autem innotescet $E$ non poterit $A$ ignorari propter datam inter eas radices differentiam præstat siquidem \ill{} $A$ per longitudinem $B$: quare factum est quod oportuit.
\marginal{$E$ ipsi}
\note{« $E$ ipsi », encadré dans la marge de gauche entre deux traits, se rapporte à « siquidem \ill{} $A$ », où un mot bref est mal formé : « $E$ præstat ipsi $A$ ». Dans « ab $E - B$ », le trait du moins est doublé.}

\subsection*{Exemplum in Cubo.}

proponatur $A$ cubus plus $B$ in $A$ quadratum ter plus $D$ plano in $A$ æquari $Z$ solido. quoniam igitur in magnitudinum climacticarum ordine, quadratum cubo proxime succedit proponitur autem hic cubus adfectus utique sub quadrato omnino æqualitas plasmatica est, alicunde effecta.
\marginal{\emph{ne}/}
\note{Dans la marge de gauche, en regard de la dernière ligne, « ne » souligné, suivi d'une barre oblique : un renvoi dont l'appel n'est pas visible dans la ligne. La phrase continue au verso, vue 72.}

\page{72}
\note{Verso du feuillet « 33 » ; l'exemple du cube, commencé au bas de la vue 71, continue. En haut à gauche, un « 33 » pâle vu par transparence. Dans la marge de gauche, des traces d'écriture pâles et inversées, décharge ou transparence, non transcrites ; les notes marginales lisibles sont données à leur place. Le bas de la page est blanc.}

est autem adfectio illa adfirmata, quare plasma fuit per additionem trientis $B$ coefficientis prout conditio cubi quæ \struck{pot} potestas est rationis triplicatæ exposcit. ad tollendum igitur plasma fiat expurgatio diatritemorion et idcirco $A + B$ esto $E$ ergo $E - B$ valebit $A$ et consequenter cubus ab $E - B$ cubus adiunctus solido a $B$ ter in \struck{\ill{}} quadratum quoque ab $E - B$ quadratum ac denique adiunctus solido a $D$ plano in $E - B$ æquabitur $Z$ solido per ea quæ proponuntur. Æqualitas igitur de $E$ secundum artem ordinetur omnibus rite peractis deprehendetur
\[ \left. \left. \begin{array}{l} E \text{ cubus} \\ +\,D \text{ plano} \\ \underline{\overline{+}}\,B \text{ quadrato} \end{array} \right\} \text{in } E \right\} \text{æquari} \ \begin{array}{l} Z \text{ solido} \\ \pm\,D \text{ plano in } B \\ -\,B \text{ cubo bis} \end{array} \]
\marginal{$\overline{\underline{\pm}}$ ×. \quad $-\,B$ quadrato ter.}
\note{Le signe devant « $B$ quadrato » est un « $+$ » surchargé entre deux traits, et la marge le corrige en « $-\,B$ quadrato ter », appelé par une croix. Devant « $D$ plano in $B$ », le signe est noirci, lu $\pm$ sans assurance. Le calcul donne $E^3 + (D - 3B^2)E = Z + DB - 2B^3$.}

quæ quidem nova æquatio pura est ab adfectione sub quadrato, qua primo proposita obruebatur cum autem innotescet $E$ non poterit $A$ ignorari, propter datam inter eas radices differentiam præstat siquidem $E$ ipsi $A$ per longitudinem $B$. quare factum est quod oportuit.

quod si \emph{quis} proponitur æqualitas potestati adfectæ per negationem directo ut cum $A$ quadratum minus $B$ in $A$ bis æquatur $Z$ plano vel $A$ cubus minus $B$ in $A$ quadratum ter adæquatur $Z$ solido $A - B$ statuitur esse $E$ et æqualitas de $E$ expositis vestigiis ordinabitur.
\marginal{\struck{\ill{}} \uncertain{et}/}
\note{« quis » est souligné ; dans la marge, un mot barré, puis « \uncertain{et} » suivi d'une barre oblique.}

Si denique proponitur æqualitas potestati adfectæ per negationem inverse ut cum $B$ in $A$ bis minus $A$ quadrato adæquatur $Z$ plano vel $B$ in $A$ quadratum ter minus $A$ cubo adæquatur $Z$ solido. $B - A$ statuitur esse $E$ et æqualitas de $E$ \emph{similibus} expositis vestigiis ordinabitur.
\marginal{\emph{similiter}.}
\note{« similibus » est souligné, et la marge porte « similiter », souligné aussi.}

Sic quadrata omnia adfecta revocantur ad pura adfecto qualitercumque, cubi ad cubos adfectos tantum sub latere adfecta qualitercumque, quadratoquadrata ad quadratoquadrata adfecta duntaxat sub latere et quadrato. adfecti qualitercumque quadratocubi ad quadratocubos adfectos duntaxat sub latere quadrato et cubo et eo deinceps progressu.

Credebant autem antiqui quoniam hac reductione æquationes quadraticas omnino expurgabant et facile explicabant
\note{La phrase s'arrête ici, sans point, au bas du texte ; elle ne se poursuit pas sur cette page.}

\page{73}
\note{Un feuillet numéroté « 34 » à l'encre en haut à droite. Au-dessus de lui, à l'angle d'un feuillet plus grand qui dépasse en haut, un « 37 » de la même encre (le feuillet de la vue 75 et suivantes). Sous le bas de la page dépasse la dernière ligne du feuillet « 35 », « $-\,D$ plano in $A$ » (vue 75), déjà visible à la vue 71 ; non transcrite ici. À gauche, les bords des feuillets précédents, avec des fins de lignes, non transcrites. La première ligne ne reprend pas la dernière phrase de la vue 72 (« Credebant autem antiqui … explicabant ») : le lien entre les deux pages n'est pas assuré.}

\uncertain{vel duobus} quoque climacticis \uncertain{accedere} ut omnino expurgarentur et a canonica puram resolutione informandum \uncertain{mechanicorum} deinde tentarunt adeo obstinate, ut aliunde methodum explicandi æquationes adfectas non \uncertain{exquisierint}.

Itaque \uncertain{erraverunt} \ill{} frustra et bonas horas mathematicas quam \uncertain{collocabant} dispendio absumpserunt. In summa methodus illa explicandi æquationes adfectas quadraticas non est catholica. Catholica quidem est methodus expurgandarum æquationum adfectione singulari salva numerorum symmetria sed non adfectione omni, ut deinceps veterum pertinaciæ non sit in \uncertain{hæredium}.
\note{Le premier mot de la page, en partie effacé, est lu sans assurance. Dans la marge de droite, en regard des trois premières lignes, les débuts de lignes d'un autre feuillet (« 1 », « rans »), non transcrits.}

Juvat autem de singulis istiusmodi reductionis formulis singula concepta theoremata, et ea in artem et usum proferre, qualia sunt quæ sequuntur.

\subsection*{De reductione quadratorum adfectorum ad pura. formulæ hæ.}

Si $A$ quadratum plus $B$ bis in $A$ æquetur $Z$ plano

$A + B$ esto $E$.

\struck{Igitur} $E$ quadratum æquabitur $Z$ plano plus $B$ quadrato

Consectarium

Itaque $L \left\{ \begin{array}{l} Z \text{ plani} \\ +\,B \text{ quadrato} \end{array} \right\} - B$ fit $A$ de qua primo quærebatur
\note{« $L$ » : \emph{latus}, la racine carrée de l'agrégat.}

Sit $B$ 1. $Z$ plano 20. $A$ $1N$.

$1Q + 2N$. æquatur 20 et fit $1N$. $L21 - 1$.

Si $A$ quadratum minus $B$ bis in $A$ æquetur $Z$ plano
\note{Au-dessus de la ligne, un « II » : le numéro de la formule, comme le « I » écrit au-dessus d'« æquetur » dans la première et le « III » en tête de la vue 74.}

$A - B$ esto $E$

\struck{Igitur} \emph{$B$} quadratum æquabitur $Z$ plano plus $B$ quadrato
\marginal{$E$}
\note{Le « $B$ » souligné est corrigé en « $E$ » dans la marge, entre deux traits.}

Consectarium

Itaque $L \left\{ \begin{array}{l} Z \text{ plani} \\ +\,B \text{ quadrato} \end{array} \right\} + B$ fit $A$ de qua primum quærebatur

Sit $B$ 1. $Z$ plano 20 $A$ $1N$.

$1Q - 2N$. æquatur 20 et fit $1N$. $L21 + 1$
\note{Le bas de la page est blanc.}

\page{74}
\note{Verso du feuillet « 34 », sans numéro ; la suite des formules de la vue 73 continue, numérotée « III » au-dessus de la ligne. À gauche, des traces pâles d'écriture, décharge ou transparence, non transcrites. Le bas de la page est blanc.}

Si $D$ bis in $A$ minus $A$ quadrato æquetur $Z$ plano

$D - E$ vel $D + E$ esto $A$

$E$ quadratum æquabitur $D$ quadrato \struck{minus} \add{minus} $Z$ plano
\note{« minus » est récrit au-dessus du même mot surchargé.}

Consectarium

Itaque $D$ minus plusve $L \left\{ \begin{array}{l} D \text{ quadrati} \\ -\,Z \text{ plano} \end{array} \right\}$ fit $A$ de qua primo quærebatur

Sit $D$ 5. $Z$ plano 20. $A$ $1N$.
\marginal{$Z$/}
\note{Dans la marge, en regard, « $Z$ » entre deux traits suivi d'une barre oblique : sans correction visible dans la ligne.}

$10N - 1Q$ æquatur 20 et fit $1N$. $5 - L5$ vel $5 + L5$.

\subsection*{De reductione cuborum simpliciter adfectorum sub quadrato ad cubos simpliciter adfectos sub latere formulæ tres.}

\[ \left. \begin{array}{l} \text{Si } A \text{ cubus} \\ +\,B \text{ ter in } A \text{ quadratum} \end{array} \right\} \text{æquetur } Z \text{ solido.} \]

$A + B$ esto $E$
\[ \left. \begin{array}{l} E \text{ cubus} \\ -\,B \text{ quadrato ter in } E \end{array} \right\} \text{æquabitur} \left\{ \begin{array}{l} Z \text{ solido} \\ -\,B \text{ cubo bis} \end{array} \right. \]

$1C + 6Q$. æquatur \uncertain{1600} est $1N$. 10

$1C - 12N$ æquatur $\underline{1684}$ est $1N$ 12
\marginal{\uncertain{fit} $1N$ 2/ \quad $\underline{1584}$}
\note{« 1684 » est souligné et corrigé en « 1584 » dans la marge, souligné aussi : $12^3 - 12 \cdot 12 = 1584 = 1600 - 2 \cdot 2^3$. Le premier nombre, écrit contre « æquatur », se lit 1600 d'après le calcul ($10^3 + 6 \cdot 10^2$), le 1 initial étant mal séparé. L'autre note de la marge, « fit $1N$ 2 » suivi d'une barre oblique, semble donner la valeur de $B$, que le texte omet.}

ad Arithmetica non incongrue \uncertain{σημεῖον} aliquod superimponitur notis alteratæ radicis ad differentiam notarum eius de qua primum quærebatur

\[ \left. \begin{array}{l} \text{Si } A \text{ cubus} \\ -\,B \text{ ter in } A \text{ quadratum} \end{array} \right\} \text{æquetur } Z \text{ solido} \]
\note{Au-dessus, le numéro de la formule, « II ».}

$A - B$ esto $E$
\[ \left. \begin{array}{l} E \text{ cubus} \\ -\,B \text{ quadrato ter in } E \end{array} \right\} \text{æquabitur} \left\{ \begin{array}{l} Z \text{ solido} \\ +\,B \text{ cubo bis} \end{array} \right. \]

$1C - 6Q$ æquatur 400 est $1N$. 10.

$1C - 12N$. æquatur 416 est $1N$ 8.

\[ \left. \begin{array}{l} \text{Si } B \text{ ter in } A \text{ quadratum} \\ -\,A \text{ cubo} \end{array} \right\} \text{æquetur } Z \text{ solido} \]
\note{Au-dessus, « III ».}

$A - B$ esto $E$.
\note{La formule s'arrête là, au bas du texte de la page.}

\page{75}
\note{Un feuillet numéroté « 35 » à l'encre en haut à droite, posé sur un feuillet plus grand numéroté « 37 », dont trois lignes dépassent au-dessus : « nempe sunt numeri triangulares 3. 6. 10. 15 \&c. proponatur autem $A$ quadratoquadratum adfectum sub $A$ et $B$ cubo simul ad expurgationem illius adfectionis quadrans » — lignes lues ici sans assurance et non transcrites comme texte. Le feuillet « 35 » est plus long que les feuillets « 33 » et « 34 » : c'est sa dernière ligne, « $-\,D$ plano in $A$ », qui dépassait sous eux aux vues 71 à 74. Il continue la formule III de la vue 74.}

\[ \left. \begin{array}{l} B \text{ quadratum ter in } E \\ -\,E \text{ cubo} \end{array} \right\} \text{æquabitur} \ \begin{array}{l} Z \text{ solido} \\ -\,B \text{ cubo bis.} \end{array} \]

vel $B - A$ esto $E$
\[ \left. \begin{array}{l} B \text{ quadratum ter in } E \\ -\,E \text{ cubo} \end{array} \right\} \text{æquabitur} \ \begin{array}{l} B \text{ cubo bis} \\ -\,Z \text{ solido} \end{array} \]

$21Q - 1C$ æquatur 972 est $1N$ 9 vel 18

$147N - 1C$ æquatur 286 est $1N$ 2 vel 11

$9Q - 1C$ æquatur 28 est $1N$. 2

$27N - 1C$. æquatur 26 est $1N$ 1.

\subsection*{De reductione cuborum adfectorum tam sub quadrato quam latere ad cubos adfectos simpliciter sub latere. formulæ septem.}

\[ \left. \begin{array}{l} \text{Si } A \text{ cubus} \\ +\,B \text{ ter in } A \text{ quadratum} \\ +\,D \text{ plano in } A \end{array} \right\} \text{æquetur } Z \text{ solido} \]
\note{Au-dessus, le numéro de la formule, « I ».}

$A + B$ esto $E$
\[ \left. \left. \begin{array}{l} E \text{ cubus} \\ +\,D \text{ plano} \\ = B \text{ quadrato ter} \end{array} \right\} \text{in } E \right\} \text{æquabitur} \left\{ \begin{array}{l} Z \text{ solido} \\ +\,D \text{ plano in } B \\ = B \text{ cubo bis.} \end{array} \right. \]
\note{Le « $=$ » est ici encore le signe de la différence, d'ordre non spécifié, et non l'égalité. Entre les deux accolades, une petite croix, sans renvoi visible.}

$1C + 30Q +$ \add{3}$30N$. æquatur 788 est $1N$ 2.
\marginal{$\underline{330}$}
\note{« 30 » est souligné, avec un « 3 » petit écrit au-dessus, en tête ; la marge donne « 330 », souligné, qui est juste : $8 + 120 + 660 = 788$.}

$1C + 30N$. æquatur 2088 est $1N$ 12.

$1C +$ \uncertain{24}$Q + 132N$ æquatur 368 est $1N$ 2

$1C - 60N$ æquatur \uncertain{400} est $1N$ 10
\note{Les deux nombres douteux sont lus d'après le calcul : avec $B = 8$ et $D = 132$, $8 + 96 + 264 = 368$ et $1000 - 600 = 400$. Le second chiffre de « 24 » ressemble à un 7, le dernier de « 400 » à un 6.}

$1C + 30Q + 4N$. æquatur 1320 est $1N$ 6.

$296N - 1C$ æquatur 640 est $1N$. 16.

\[ \left. \begin{array}{l} \text{Si } A \text{ cubus} \\ +\,B \text{ ter in } A \text{ quadratum} \\ -\,D \text{ plano in } A \end{array} \right\} \text{æquetur } Z \text{ solido} \]
\note{Au-dessus, « II ». La formule s'arrête là ; le bas du feuillet est blanc.}

\page{76}
\note{Verso du feuillet « 35 », sans numéro ; la formule II, dont l'énoncé termine la vue 75, se poursuit. Tout est de la même main, en formules et en exemples numériques, sans prose. Le bas de la page est blanc.}

$A + B$ esto $E$
\[ \left. \left. \begin{array}{l} E \text{ cubus} \\ = B \text{ quadrato ter} \\ = D \text{ plano} \end{array} \right\} \text{in } E \right\} \text{æquabitur} \ \begin{array}{l} Z \text{ solido} \\ = B \text{ cubo bis} \\ = D \text{ plano in } B \end{array} \]

$1C + 6Q - 48N$. æquatur $\underline{512}$ est $1N$ 8.
\marginal{\struck{\ill{}}}
\note{« 512 » est souligné ; dans la marge, entre deux traits, un nombre barré, \ill{}. Le calcul confirme 512 : $512 + 384 - 384$.}

$1C -$ \uncertain{60}$N$ æquatur 400 est $1N$ 10

$1C + 30Q - 48N$ æquatur 32 est $1N$ 2

$348N - 1C$ æquatur 2448 est $1N$ 12.
\note{Le coefficient de la deuxième ligne est surchargé ; 60 est ce que le calcul demande ($1000 - 600 = 400$).}

\[ \left. \begin{array}{l} \text{Si } A \text{ cubus} \\ -\,B \text{ ter in } A \text{ quadratum} \\ = D \text{ plano in } A \end{array} \right\} \text{æquetur } Z \text{ solido} \]
\marginal{$\overline{\underline{+}}$}
\note{Au-dessus, le numéro « III ». Le signe devant « $D$ plano in $A$ » est un « $=$ » repassé ; la marge porte « $+$ » entre deux traits.}

$A - B$ esto $E$
\[ \left. \left. \begin{array}{l} E \text{ cubus} \\ = B \text{ quadrato ter} \\ +\,D \text{ plano} \end{array} \right\} \text{in } E \right\} \text{æquabitur} \ \begin{array}{l} Z \text{ solido} \\ +\,B \text{ cubo bis} \\ = D \text{ plano in } B. \end{array} \]
\note{Au-dessus de « ter », un signe « $\mp$ » ajouté.}

$1C - 30Q + 330N$ æquatur 1368 est $1N$ 12

$1C + 30N$. æquatur 68 est $1N$ 2.

$1C$ \ill{} $12Q + 28N$. æquatur $\underline{76}$ est $1N$ 10
\marginal{80}
\note{Le signe après « $1C$ » est noyé dans une tache ; le calcul demande $-$. « 76 » est souligné et corrigé en « 80 » dans la marge : $1000 - 1200 + 280 = 80$.}

$1C - 20N$ æquatur 96 est $1N$ 6.

$1C - 18Q + 88N$ æquatur 80 est $1N$ 10

$20N - 1C$ æquatur 16 est $1N$ 4.

vel $B - A$ esto $E$
\[ \left. \begin{array}{l} B \text{ quadrato ter in } E \\ = D \text{ plano} \\ = E \text{ cubo} \end{array} \right\} \text{æquabitur} \ \begin{array}{l} Z \text{ solido} \\ +\,B \text{ cubo bis} \\ = D \text{ plano in } B \end{array} \]

$1C - 30Q + 200N$ æquatur 336 est $1N$ 6

$100N - 1C$ æquatur 336 est $1N$ 4

$1C - 30Q + 280N$ æquatur 704 est $1N$ 4

$1C - 20N$ æquatur 96 est $1N$ 6

$1C - 30Q + 330X$ æquatur 1232 est $1N$ 8
\marginal{$\overline{\underline{N}}$.}
\note{« $330X$ » : le $N$ mal formé, corrigé en « $N$ » dans la marge.}

$1C + 30N$ æquatur 68 est $1N$ 2.

\page{77}
\note{Un feuillet numéroté « 36 » à l'encre en haut à droite, posé sur le grand feuillet « 37 », dont on voit au-dessus les mêmes trois lignes qu'à la vue 75 (« nempe sunt numeri triangulares 3. 6. 10. 15 \&c. … quadrans »), et à droite une bande avec des fins de lignes ; non transcrites ici. Le feuillet « 36 » continue les sept formules de la vue 75, numérotées « IIII » à « VI » au-dessus de chaque énoncé. Le bas de la page est blanc.}

\[ \left. \begin{array}{l} \text{Si } A \text{ cubus} \\ -\,B \text{ ter in } A \text{ quadratum} \\ -\,D \text{ plano in } A \end{array} \right\} \text{æquetur } Z \text{ solido} \]
\note{Au-dessus, « IIII ».}

$A - B$ esto $E$
\[ \left. \left. \begin{array}{l} E \text{ cubus} \\ -\,B \text{ quadrato ter} \\ -\,D \text{ plano} \end{array} \right\} \text{in } E \right\} \text{æquabitur} \left\{ \begin{array}{l} Z \text{ solido} \\ \underline{\phantom{+}}\,B \text{ cubo bis} \\ \underline{\phantom{+}}\,D \text{ plano in } B \end{array} \right. \]
\marginal{$\overline{\underline{+}}$ \quad $\overline{\underline{+}}$}
\note{Devant « $B$ cubo bis » et « $D$ plano in $B$ », le signe est laissé en blanc sur un trait ; la marge donne les deux « $+$ », entre des traits. Le signe devant « $B$ quadrato ter » est peu net.}

$1C - 6Q - 28N$. æquatur $\underline{\overline{20}}$ est $1N$ 10
\marginal{120.}
\note{« 20 » est souligné et surligné, et corrigé en « 120 » dans la marge : $1000 - 600 - 280 = 120$.}

$1C - 40N$ æquatur 192 est $1N$ 8.

\[ \left. \begin{array}{l} \text{Si } D \text{ plano in } A \\ -\,B \text{ ter in } A \text{ quadratum} \\ -\,A \text{ cubo} \end{array} \right\} \text{æquatur } Z \text{ solido} \]
\note{Au-dessus, « V ». Après le second « $B$ », une lettre barrée.}

$A + B$ esto $E$
\[ \left. \left. \begin{array}{l} D \text{ plano} \\ +\,B \text{ quadrato ter} \end{array} \right\} \text{in } E \ \begin{array}{l} \\ -\,E \text{ cubo} \end{array} \right\} \text{æquabitur} \ \begin{array}{l} Z \text{ solido} \\ \underline{\phantom{+}}\,B \text{ cubo bis} \\ \underline{\phantom{+}}\,D \text{ plano in } B \end{array} \]
\marginal{$\overline{\underline{+}}$ \quad $\overline{\underline{+}}$}
\note{Mêmes signes laissés en blanc, et mêmes « $+$ » dans la marge : $400 \cdot 12 - 12^3 = 72 + 2000 + 1000 = 3072$.}

$100N - 30Q - 1C$ æquatur 72 est $1N$ 2

$400N - 1C$ æquatur 3072 est $1N$ 12.

\[ \left. \begin{array}{l} \text{Si } B \text{ ter in } A \text{ quadratum} \\ +\,D \text{ plano in } A \\ -\,A \text{ cubo} \end{array} \right\} \text{æquetur } Z \text{ solido} \]
\note{Au-dessus, « VI ».}

$A - B$ esto $E$
\[ \left. \left. \begin{array}{l} D \text{ plano} \\ +\,B \text{ quadrato ter} \end{array} \right\} \text{in } E \ \begin{array}{l} \\ -\,E \text{ cubo} \end{array} \right\} \text{æquabitur} \ \begin{array}{l} Z \text{ solido} \\ -\,D \text{ plano in } B \\ -\,B \text{ cubo bis} \end{array} \]

$18Q + 92N - 1C$ æquatur 1720 est $1N$ 10

$200N - 1C$ æquatur 736 est $1N$ 4

vel $B - A$ esto $E$
\[ \left. \left. \begin{array}{l} D \text{ plano} \\ +\,B \text{ quadrato ter} \end{array} \right\} \text{in } E \ \begin{array}{l} \\ -\,E \text{ cubo} \end{array} \right\} \text{æquabitur} \left\{ \begin{array}{l} B \text{ cubo bis} \\ \underline{\phantom{+}}\,D \text{ plano in } B \\ -\,Z \text{ solido} \end{array} \right. \]
\marginal{$+$}
\note{Devant « $D$ plano in $B$ », le signe laissé en blanc sur un trait ; la marge porte « $+$ ».}

$\underline{\overline{300}}Q + 100N - 1C$ æquatur 1464 est $1N$ 6
\marginal{$\underline{\overline{30}}$}
\note{« 300 » est corrigé en « 30 » dans la marge : $30 \cdot 36 + 600 - 216 = 1464$.}

$400N - 1C$ æquatur 1536 est $1N$ 4.
\note{La septième formule annoncée à la vue 75 n'est pas sur cette page.}

\page{78}
\note{Verso du feuillet « 36 », sans numéro ; la septième formule, annoncée à la vue 75, ouvre la page, numérotée « VII » ; à côté, plus pâle, « IIII », transparence du recto. Le texte en prose qui suit revient à l'expurgation par onces. Le bas de la page est blanc.}

\[ \left. \begin{array}{l} \text{Si } B \text{ ter in } A \text{ quadratum} \\ -\,D \text{ plano in } A \\ -\,A \text{ cubo} \end{array} \right\} \text{æquetur } Z \text{ solido.} \]

$A - B$ esto $E$
\[ \left. \left. \begin{array}{l} B \text{ quadratum ter} \\ = D \text{ plano} \end{array} \right\} \text{in } E \ \begin{array}{l} \\ = E \text{ cubo} \end{array} \right\} \text{æquabitur} \ \begin{array}{l} Z \text{ solido} \\ +\,D \text{ plano in } B \\ = B \text{ cubo bis} \end{array} \]

$18Q - 78N - 1C$ æquatur 20 est $1N$ 10

$30Q - 1C$ æquatur 56 est $1N$ 4
\note{« $30Q$ » : ainsi, sans correction ; le calcul demande $30N$ ($30 \cdot 4 - 64 = 56$).}

$12Q - 18N - 1C$ æquatur 20 est $1N$ 10

$1C - 30N$ æquatur $\underline{96}$ est $1N$ 6.
\marginal{$\underline{36}$}
\note{« 96 » est souligné et corrigé en « 36 » dans la marge ; l'égalité juste, $30 \cdot 6 - 6^3 = -36$, s'écrirait $30N - 1C$ æquatur $-36$, soit $1C - 30N$ æquatur 36.}

vel $B - A$ esto $E$
\[ \left. \left. \begin{array}{l} B \text{ quadratum ter} \\ -\,D \text{ plano} \end{array} \right\} \text{in } E \ \begin{array}{l} \\ -\,E \text{ cubo} \end{array} \right\} \text{æquabitur} \left\{ \begin{array}{l} B \text{ cubo bis} \\ -\,D \text{ plano in } B \\ -\,Z \text{ solido} \end{array} \right. \]

$\underline{\overline{300}}Q - 100N - 1C$ æquatur 264 est $1N$ 6
\marginal{$\underline{30}$}
\note{« 300 » corrigé en « 30 » dans la marge, comme à la vue 77 : $1080 - 600 - 216 = 264$.}

$200N - 1C$ æquatur 736 est $1N$ 4.

\uncertain{et} autem expurgatione per uncias easque simplices liberantur æquationes adfectione sub gradu qui in climacticarum ordine potestatum proxime subsequitur. \uncertain{sic eodem} per uncias triangulares et pyramidales et ex iis compositas possunt æquationes liberari adfectione sub gradibus inferioribus, cum quis considerata coefficiente sub graduali et potestate suæque radicis tot uncias adsumendo quot × dictat genesis potestatum purarum a binomia radice ut cum coefficientis radici de qua quæritur homogenea per uncias simplices taxantur sic homogenea radicis quadrato per uncias triangulares cubo per pyramidales et eo continuo ordine.
\marginal{|o| \quad \uncertain{et}/ \quad × requirit.}
\note{Dans la marge de gauche, en regard des deux premières lignes, « |o| » et « et » suivi d'une barre oblique, sans appel visible dans la ligne ; plus bas, « requirit. », appelé par la croix écrite après « quot » : « quot requirit genesis ». Dans « climacticarum » et « potestatum », les finales sont soulignées.}

proponatur $A$ cubus adfectus sub $A$ et $B$ quadrato sumetur ad expurgationem triens $B$

et si proponatur $A$ quadratoquadratum adfectum sub $A$ quadrato et $B$ quadrato sumetur ad expurgationem illius adfectionis sextans $B$.

et si proponatur $A$ quadratocubus adfectus sub $A$ cubo et $B$ \emph{cubo} sumetur ad expurgationem decima pars $B$.
\marginal{quadrato/}
\note{« cubo » est souligné et corrigé en « quadrato » dans la marge. Les triens, sextans et « decima pars » sont les inverses des nombres triangulaires 3, 6, 10 : les lignes « nempe sunt numeri triangulares 3. 6. 10. 15 \&c. » du feuillet « 37 », vues aux vues 75 et 77, semblent la suite de cette page.}

\page{79}
\note{Le grand feuillet numéroté « 37 » à l'encre en haut à droite, dont les trois premières lignes se voyaient aux vues 75 et 77 au-dessus des feuillets « 35 » et « 36 ». Il continue le propos de la fin de la vue 78 sur les onces triangulaires. Sous le bas de la page dépasse un feuillet plus long, qui porte « metathesi $A$ cubus æquatur $B$ plano in $A$ $+\,Z$ solido » (un mot barré entre les deux lignes) : lignes d'un feuillet suivant, hors de ce lot, non transcrites ici. Dans la marge de gauche, des traces pâles, décharge, non transcrites.}

nempe sunt numeri triangulares 3. 6. 10. 15 \&c.

proponatur autem $A$ quadratoquadratum adfectum sub $A$ et $B$ cubo sumetur ad expurgationem illius adfectionis quadrans $B$

et si proponatur $A$ quadratocubus adfectus sub $A$ quadrato et $B$ cubo sumetur ad expurgationem decima pars $B$.

nempe sunt numeri pyramidales 4. 10. 20. 35 \&c.

quæ quemadmodum ad ulteriora possunt aptari satis fit manifestum, valet autem illa unciarum triangularium pyramidalium et ex iis compositarum expurgatio dum illa ea adfectione qua liberanda proponitur obruitur \emph{æqualitate} in cuius autem adfectionis deletæ locum, succedunt adfectiones singulæ sub aliis gradibus reliquis.
\marginal{tas}
\note{La finale de « æqualitate », soulignée, est corrigée en « -tas » dans la marge : « obruitur æqualitas ».}

\subsection*{De transmutatione πρῶτον ἔσχατον quæ remedium est adversus vitium negationis. Cap. II}

Æquationes in quibus homogenea \emph{æquationum} validiora negantur de potestate utiliter per transmutationem πρῶτον ἔσχατον emendantur. ea fit per analogiam rationis simplicis, \emph{explicando} homogeneorum comparationis ad ipsam radicem de qua quæritur unde oriatur incerta alia radix sub cuius specie æquationem primo propositam liceat ducere et novam ordinare.
\marginal{adfectionum. \quad applicando}
\note{Les deux mots soulignés sont corrigés dans la marge, entre des traits : « adfectionum » pour « æquationum », « applicando » pour « explicando ». Après « rationis », un mot barré ou surchargé.}

Sic enim adfectiones illæ negatæ transeunt in affirmatas et e contra salva numerorum symmetria

prodest etiam in \uncertain{vitium} ad asymmetrias

nomen autem πρῶτον ἔσχατον sortita est ab eo ad quem proposita primo æquatio \emph{vocatur} analogismo quoniam in eius formula terminus de quo primo quærebatur est primus, is qui post metamorphosin \emph{fit} postremus vel e converso.
\marginal{renovatur \quad primus, fit}
\note{« vocatur », souligné et surchargé, est corrigé en « renovatur » dans la marge ; « primus, fit », dans la marge, semble reprendre les mots « est primus, is » et « fit », ce dernier souligné dans la ligne.}

quæ quoniam ex ipsa operis formula facilius percipiuntur

proponatur $\left. \begin{array}{l} A \text{ cubus} \\ \quad -\,B \text{ plano in } A \end{array} \right\}$ æquari $Z$ solido

sit autem explicanda æqualitas

quoniam igitur $Z$ solidum potestas est affecta negate, de negatis autem ars non statuitur, quæ proponitur æquatio inexplicabilem primo transmutanda est qualis quæ adficeretur affirmate

quocirca $\frac{Z \text{ solidum}}{A}$ esto $E$ plano ergo $\frac{Z \text{ solidum}}{E \text{ plano}}$ erit $A$
\note{Le texte s'arrête là, au milieu de la page ; le reste du feuillet est blanc.}

\page{80}
\note{Verso du grand feuillet « 37 », sans numéro ; le chapitre II continue sans rupture la vue 79. Sous le bas de la page dépasse le feuillet « 36 » : on y lit la dernière ligne de la vue 78, « sumetur ad expurgationem decima pars $B$ », avec la note marginale « quadrato/ » ; non transcrite ici. La marge de gauche porte des traces pâles, décharge, et une longue note de la main du texte, plus serrée. La dernière ligne s'arrête sur « itaque ex iis quæ proponuntur » : la suite est hors de ce lot.}

unde ex iis quæ proponuntur
\[ \left. \begin{array}{l} \dfrac{Z \text{ solido solido solidum}}{E \text{ plano plano plano}} \\[2ex] \text{minus } B \text{ plano in } \dfrac{Z \text{ solidum}}{E \text{ plano}} \end{array} \right\} \text{æquabitur } Z \text{ solido} \]

ducantur omnia in $E$ planoplanoplano
\[ \left. \begin{array}{l} Z \text{ solidosolidosolidum} \\ \text{minus } B \text{ plano in } Z \text{ solidum in } E \text{ plano} \end{array} \right\} \text{æquabitur } Z \text{ solido in } E \text{ planoplanoplano} \]

et omnibus per $Z$ solidum divisis adhibitaque congrua antithesi
\[ \left. \begin{array}{l} E \text{ planoplanoplano} \\ +\,B \text{ plano in } E \text{ planoplano} \end{array} \right\} \text{æquabitur } Z \text{ solidosolido} \]

\subsection*{Aliter}

\[ \left. \begin{array}{l} E \text{ plani cubus} \\ +\,B \text{ plano in } E \text{ plani quadratum} \end{array} \right\} \text{æquabitur } Z \text{ solidi quadrato} \]

quæ æquatio omnino explicabilis est ex tradita analysi potestatum cubicarum quæ adficiuntur sub dato coefficiente adfirmate. Innotescit igitur $E$ planum ad quod cum applicabitur $Z$ solidum exurget $A$ de qua primo quærebatur

Analogismus autem æqualitatis de $A$ inveniendus erat ut $A$ ad $LC$ $Z$ solidi ita $LC$ $Z$ solidosolidi ad $A$ quadratum $-\,B$ plano. vel est ad $E$ planum

quo pertinet nominis πρῶτον ἔσχατον notatio
\note{« $LC$ » : \emph{latus cubicum}. Le « $-\,B$ plano » est écrit sous « $A$ quadratum ».}

\marginal{analogismus primæ æquationis est ut $A$ ad $Z$ \ill{} lat. $Z\,s$. ita $Z\,s$. vel \ill{} ad \ill{} $E$ pl. quoniam \ill{} est \struck{\ill{}} ut $A$ radix æquationis primo propositæ ad latus cub. homogenei comparationis $Z$ solidi ita latus cub. homogenei comparationis in æquatione secunda \ill{} est $LC$ $Z\,ss$. ad $E$ pl. radicem eiusdem æquationis secundæ \quad $+\,96$. \quad $\underline{1600}$}
\note{Longue note de la marge de gauche, en regard de l'« Aliter » et du « Analogismus », fermée par un trait ; ses quatrième et cinquième lignes sont raturées et tachées, et la lecture n'en est que partielle. Au bas, « $+\,96$ » et « 1600 » souligné, séparés d'une barre : les corrections de l'exemple qui suit.}

proponatur $1C - 96N$ æquari 40

effecta $\frac{40}{1N}$ esse $1N$. $1C \mathbin{\underline{-\,46}}Q$ æquabitur $\underline{160}$ et fit $1N$ 4

quare $\frac{40}{4}$ seu 10 est radix primo quæsita
\note{« $-\,46$ » et « 160 » sont soulignés, corrigés dans la marge en « $+\,96$ » et « 1600 » : $4^3 + 96 \cdot 4^2 = 1600 = 40^2$.}

et si accidat in proposita primo æquatione cubica homogeneum comparationis esse in sua soliditate irrationale sed facultate veluti \emph{quadrata} rationale ea asymmetria in æquatione nova evanescet \struck{et}.
\marginal{× \emph{quadratica} \quad \emph{proponatur}.}
\note{« quadrata », souligné, est corrigé en « quadratica » dans la marge ; « proponatur », souligné aussi dans la marge, est appelé par la croix devant l'exemple qui suit, dont il complète le début.}

$1C - 10N$ \emph{æquari} $L48$ effecta $\frac{L48}{1N}$ esse $1N$.

$1C + 10Q$ æquabitur 48 et fiet $1N$ 2

quare $\frac{L48}{L4}$ seu latus 12 est radix primo quæsita
\note{« $L$ » : \emph{latus}, la racine carrée. $E^3 + 10E^2 = 48$ pour $E = 2$, et $A = \sqrt{48}/2 = \sqrt{12}$.}

rursus proponatur $A$ quadratoquadratum minus $B$ in $A$ cubum minus $D$ plano in $A$ quadratum æquari $Z$ planoplano \uncertain{sitque} explicanda æqualitas. quoniam igitur $Z$ planoplanum est potestas affecta negate, de negatis autem ars non statuitur quæ proponitur æquatio inexplicabilis transmutanda est

quocirca $\frac{Z \text{ planoplanum}}{A}$ esto $E$ solidum ergo $\frac{Z \text{ planoplanum}}{E \text{ solido}}$ erit $A$ itaque ex iis quæ proponuntur

\end{document}
